\documentclass[11pt,a4paper]{article}

\usepackage[margin=2.4cm]{geometry}
\usepackage{amsmath}
\usepackage{booktabs}
\usepackage{array}
\usepackage{enumitem}
\usepackage{titlesec}
\usepackage{hyperref}
\hypersetup{colorlinks=true, linkcolor=black, urlcolor=blue}

% headings: bold, barely larger than the text, well spaced
\titleformat{\section}{\normalfont\large\bfseries}{\thesection.}{0.6em}{}
\titleformat{\subsection}{\normalfont\normalsize\bfseries}{\thesubsection.}{0.6em}{}
\titlespacing*{\section}{0pt}{1.5em}{0.5em}
\titlespacing*{\subsection}{0pt}{1.2em}{0.4em}

\setlength{\parskip}{0.45em}
\setlength{\parindent}{0pt}
\newcommand{\unit}[1]{\,\mathrm{#1}}

% long class and method names in \texttt cannot be hyphenated; let TeX loosen a
% line rather than push it into the margin
\setlength{\emergencystretch}{3em}

\begin{document}

\begin{center}
  {\large\bfseries A guide to the analyzer}\\[0.35em]
  {\normalsize How the code turns the bench readings into results}\\[0.9em]
  {\small Suresh Hemal \quad\textbullet\quad 2 October 2026}
\end{center}
\vspace{0.4em}
\hrule
\vspace{1.2em}

\tableofcontents
\vspace{1.2em}

This document describes exactly what the analyzer does to the numbers, stage by stage, so that
any figure the report prints can be checked by hand. It names the real classes and methods, so
each description can be read against the source beside it, and it says why each choice was made
wherever a choice was made --- which readings are set aside and on what grounds, how every
quantity is formed from the raw weights and lab concentrations, what one whole run looks like
worked through from balance reading to permselectivity, how the polynomial is fitted and chosen,
and what the present set of runs can and cannot be asked.

\section{What the analyzer does, end to end}

The analyzer turns four CSV files into two things: one row of numbers per run, and a fitted
polynomial per membrane with the best setup that polynomial allows. Nothing in between is
saved to disk. Every stage takes objects in and hands objects on, and the whole chain is
three statements long.

\subsection{The four stages, in the order they run}

\texttt{Analyzer.analyze} in \texttt{src/limodel/analyzer/analyzer.py} is the order, and it is
short on purpose --- the order is the only thing this class knows:

\begin{enumerate}[nosep]
  \item \textbf{Load.} \texttt{DataLoader.load} reads the four CSV files and returns a list of
        \texttt{Run} objects, each with its readings already inside it.
  \item \textbf{Preprocess.} \texttt{Preprocessor.run} takes that list and gives back the
        shorter list the modelling may use. On the way it clears single values it cannot
        believe.
  \item \textbf{Calculate.} \texttt{Calculator.find\_results} takes the surviving runs
        and returns one \texttt{Result} per run whose two outcomes came out.
  \item \textbf{Model.} \texttt{Surrogate.find\_models} takes those results, groups them by
        membrane, fits three outcomes per membrane, and searches each membrane for its best
        setup.
\end{enumerate}

Each stage sees only what the one before it returned. The preprocessor never sees a CSV file;
the calculator never sees a run that was dropped; the fitting never sees a \texttt{Run} at all,
only \texttt{Result} rows.

What comes back is an \texttt{Analysis}: the runs used, the results, and the models. The runs
travel beside the results rather than being replaced by them, because a reader checking a
number wants the readings it came from, and because the count of runs used is part of the
claim --- a polynomial from sixteen runs and one from twenty-one are different statements.

\subsection{The journal: made once, handed down}

The \texttt{Journal} is created empty in one place --- the \texttt{analyze} command in
\texttt{cli/main.py} --- and passed into \texttt{Analyze}, which passes it into
\texttt{Analyzer}, which passes it into \texttt{Preprocessor}, \texttt{Calculator} and
\texttt{Surrogate}. There is exactly one of them per run of the program.

It is injected rather than owned by each stage for one reason: the report has to add up. If the
preprocessor kept its own list of drops and the calculator kept another, the two would have to
be collected and merged by whoever printed them, and a stage added later would be easy to leave
out of that merge. With one journal, a stage that writes nothing down has simply written nothing
down, and the arithmetic at the top of the report is a count of one list.

Two things get written. A \textbf{drop} means the whole record is out of the modelling. A
\textbf{skip} means the record stays and one of its values does not. An \texttt{Entry} carries
the action, the record's id, the record's class name, the reason in plain words, and (for a
skip) the field that was cleared.

The journal takes the id and the class name off the record itself, inside
\texttt{Journal.\_write}, instead of having each caller pass them. That is deliberate: a caller
that passed its own id could pass one that disagreed with the record, and nothing would catch
it.

Because every removal leaves an entry, the line under the journal reconciles by construction:

\begin{equation}
  \label{eq:reconcile}
  N_{\text{loaded}} \;=\; N_{\text{used}} \;+\; \bigl|\{\, \text{drop ids}\,\} \setminus
  \{\, \text{used ids}\,\}\bigr|
\end{equation}

The set difference in \eqref{eq:reconcile} matters. A run can be dropped by the calculator
\emph{after} the preprocessor already passed it, so it would appear both in the used list
handed over and in the drops; subtracting the used ids counts each run once. On the current
data the analyzer writes 32 entries --- 5 drops and 27 skips --- and the report prints
\texttt{using 16 of 21 runs}.

The \texttt{Surrogate} stage writes nothing to the journal. A membrane it will not fit is
reported on its own answer, in \texttt{SurrogateModels.not\_fitted}, because that is a finding
about the design of the experiments rather than a record that was set aside.

\subsection{The four files, and why loading is all or nothing}

\begin{table}[htbp]
  \centering
  \caption{The four CSV files in \texttt{data/}, with the row counts in the current data.}
  \label{tab:files}
  \begin{tabular}{@{}llr@{}}
    \toprule
    File & One row per & Rows \\
    \midrule
    \texttt{runs.csv}           & experiment, with how it was set up        & 21 \\
    \texttt{tank\_readings.csv} & run, minute and tank                      & 1049 \\
    \texttt{stack\_readings.csv}& run and minute, from the power log        & 891 \\
    \texttt{samples.csv}        & bottle sent to the lab                    & 345 \\
    \bottomrule
  \end{tabular}
\end{table}

\texttt{runs.csv} holds the setup and nothing measured during the run: the membranes, the feed
concentrations, the set current, the flow, the membrane area, the cell pairs, the three tank
tares, and the \texttt{include} flag. \texttt{tank\_readings.csv} holds conductivity, pH,
temperature and weight for the concentrate, dilute and rinse.
\texttt{stack\_readings.csv} holds the two voltage columns and the two current columns ---
supply and multimeter --- which is what lets the checks compare one instrument against the
other. \texttt{samples.csv} holds the lab's answers: lithium, magnesium, sodium, potassium,
calcium and TOC.

Each file has a schema in \texttt{src/limodel/schema/}, written as a \emph{pandera}
\texttt{DataFrameModel}. Every schema sets \texttt{coerce = True} and \texttt{strict = True}, so
values are converted to the declared type on load and the columns must be exactly the declared
ones --- no extras, none missing. The three reading files declare a uniqueness key
(\texttt{run\_code} with \texttt{minute}, plus \texttt{tank} where there is one), and
\texttt{run\_code} and \texttt{sample\_id} are unique in their own files. Measured columns are
\texttt{nullable}; setup columns are not, and most carry a range (\texttt{li\_feed\_mg\_l}
greater than zero, \texttt{ph} between 0 and 14, \texttt{area\_cm2} greater than zero).

A blank cell is missing, never zero. \texttt{pandas} reads it as NaN, and the loader's
\texttt{\_without\_missing} turns it into \texttt{None} before any model is built, so a value
that was not measured never arrives as a number that was.

Loading is all or nothing. Every accessor on \texttt{CsvSource} is decorated with
\texttt{@pa.check\_types}, so the table is validated as it is handed out; one offending cell
raises, and no \texttt{Run} object is built at all. That is the right behaviour here because the
alternative is worse: a table loaded with the bad rows quietly dropped would still produce a
polynomial, and it would be the wrong polynomial with no sign that anything had gone missing.
Fixing data is a separate job, which is why \texttt{CsvSource.raw\_tables} exists --- it returns
the tables unchecked, and the \texttt{check} command uses it to list what to repair in the
workbook. Calculating refuses bad data; reporting on bad data does not.

\subsection{The four objects, and how they nest}

\begin{itemize}[nosep]
  \item \texttt{Run} --- the setup, plus two lists: \texttt{stack\_readings} and
        \texttt{tank\_readings}.
  \item \texttt{StackReading} --- the power log at one minute of a run.
  \item \texttt{TankReading} --- one tank at one minute of a run, with an optional
        \texttt{sample}.
  \item \texttt{Sample} --- the lab's results for liquid drawn from one tank at one minute.
\end{itemize}

All four are \emph{pydantic} models, so the nesting is also the validation: handing
\texttt{Run.model\_validate} a dictionary of rows builds and checks the whole tree in one call.

A sample hangs off the tank reading it was drawn from, not off the run. The reason is that a
sample is not a property of a run; it is a property of a moment in a tank. Two tanks are often
sampled at the same minute, so a flat list on the run would need the tank and the minute carried
along to say which bottle belonged where --- and that pairing would then exist twice, once on
the sample and once in whatever lookup used it. Hung off the reading, the pairing is the
structure. \texttt{DataLoader.\_with\_samples} does the matching once, on
\texttt{(tank, minute)}, and after that any code that has a reading has its sample and its
conductivity, pH, temperature and weight together. The end-sample step depends on exactly this:
it walks one tank's readings, keeps the ones that have a sample, sorts them by minute, and looks
at the first and the last.

On the current data, 21 runs carry 1049 tank readings and 891 stack readings, and all 345
samples attach to a reading.

\subsection{The id on every object}

Every object carries an \texttt{id}, and the ids are built so that one string says where a
value came from.

\begin{table}[htbp]
  \centering
  \caption{What each object's \texttt{id} is, with an example from the current data.}
  \label{tab:ids}
  \begin{tabular}{@{}lll@{}}
    \toprule
    Object & \texttt{id} & Example \\
    \midrule
    \texttt{Run}          & the run code                          & \texttt{AJ} \\
    \texttt{TankReading}  & run, tank, minute                     & \texttt{AJ:conc:0} \\
    \texttt{StackReading} & run, minute                           & \texttt{AJ:0} \\
    \texttt{Sample}       & the lab's own label                   & \texttt{AE1} \\
    \texttt{Result}       & the run code                          & \texttt{AE} \\
    \bottomrule
  \end{tabular}
\end{table}

A \texttt{Sample} keeps the lab's label as its id, because that is the string written on the
bottle and the one to quote when querying a result; it carries \texttt{run\_id} and
\texttt{tank\_id} as well, so it can still point back at the reading it came from.
\texttt{TankReading} and \texttt{StackReading} each carry \texttt{run\_id} for the same reason.
\texttt{Result} has an id at all only so that the journal can accept one.

That last point is what the ids are for. \texttt{Journal.\_write} does not care what it has been
handed; it reads \texttt{record.id} and \texttt{type(record).\_\_name\_\_} and writes those down.
So any object with an id can be journalled, and the entry locates the thing without the journal
knowing anything about it. The report then splits the id on colons to sort the entries by run
and then by minute, which is why a dropped run (no minute in its id) prints above its own
readings.

\subsection{The stages and the classes that do them}

\begin{table}[htbp]
  \centering
  \small
  \caption{Each stage, the class that drives it, and what it hands on.}
  \label{tab:stages}
  \begin{tabular}{@{}lll@{}}
    \toprule
    Stage & Class & Hands on \\
    \midrule
    Load       & \texttt{DataLoader}, \texttt{CsvSource}      & \texttt{list[Run]} \\
    Preprocess & \texttt{Preprocessor} and three \texttt{Step}s & \texttt{list[Run]}, shorter \\
    Calculate  & \texttt{Calculator} and the quantity classes & \texttt{list[Result]} \\
    Model      & \texttt{Surrogate}, \texttt{FittedOutcome}, \texttt{DesignPoints}
               & \texttt{SurrogateModels} \\
    \midrule
    Order      & \texttt{Analyzer}  & \texttt{Analysis} \\
    Record     & \texttt{Journal}   & \texttt{tuple[Entry, ...]} \\
    Report     & \texttt{Analyze}   & the printed lines \\
    \bottomrule
  \end{tabular}
\end{table}

\subsection{What follows}

The rest of this document opens each stage in turn: the three preprocessing steps and the six
reading checks, with the thresholds each one judges by and the reason behind every number; the
quantities the calculator works out from tank weights, lab results and the power log, and the
assumptions inside each one; the second-order polynomial, how the two settings are logged and
coded, how terms are kept or dropped, and how leave-one-out error is weighed against the bench's
own scatter; and the search for the best setup under the recovery floor and the voltage ceiling.

\section{Deciding which readings can be believed}

Between loading the CSVs and working out any chemistry sits one class, \texttt{Preprocessor}.
It takes the list of \texttt{Run} objects the loader built and returns a shorter list, with
some values inside the survivors set to \texttt{None}. It saves nothing and prints nothing.
Everything it removes is written to the \texttt{Journal} it was handed, so the report can say
later exactly what the modelling never saw.

Two things it deliberately does not do. It does not repair a reading: there is no filling in,
no smoothing, no interpolation. A value it cannot believe is removed, never replaced. And it
does not check that required fields are present. That is the loader's job and the loader is
all-or-nothing: each CSV is validated against its pandera schema as it is handed out, and a
table with a required field missing raises rather than loading. So by the time the
preprocessing sees a run, the run is well formed and only its values are in question. There is
a comment in \texttt{Preprocessor.steps} saying exactly this, marking where a fourth step used
to be needed and no longer is.

\subsection{The three steps, in order}

\texttt{Preprocessor.steps} is a tuple of three \texttt{Step} objects. Each one is given the
runs the step before it returned, so the order is the order of the pipeline.

\begin{enumerate}[nosep]
  \item \textbf{\texttt{DropExcludedRuns}} --- keeps only runs with \texttt{include} true.
        The mark is a decision taken at the bench, and this step only obeys it. If it would
        return nothing while it was given something, it raises \texttt{ValueError} instead:
        an empty result means the wrong file was loaded, not that every run was bad. Its
        journal label is \emph{not marked for use}.
  \item \textbf{\texttt{SkipSuspectReadings}} --- drops no run. It clears single values it
        cannot believe, and writes down why. Its label is \emph{suspect reading}.
  \item \textbf{\texttt{DropRunsMissingAnEndSample}} --- keeps only runs the lab has answered
        for: both lithium and magnesium, in both the concentrate and the dilute, at the
        first sampled minute and the last. Its label is \emph{no lab results at both ends}.
\end{enumerate}

Why step 3 comes last rather than first: it is cheaper to drop a run before inspecting its
readings, but then a run on its way out would never be looked over. Running the inspection
first means a doomed run still gets its bad values journalled, which is how the bench learns
that a particular sheet has a weight written down wrong even though that run never reached the
model. The real output shows this working --- run \texttt{AU} is dropped, and its rinse weight
at minute 240 is journalled anyway.

Step 3 asks only for the two ends because that is all the chemistry needs: the amount that
crossed the membrane is a difference between a start and a finish. The samples in between are
not asked for. The rinse is not asked for at all, because it is not part of the chemistry the
modelling uses. Both metals are required rather than lithium alone, because
permselectivity is a ratio of ratios and one metal cannot give it.

\subsection{Steps hold inspectors, inspectors hold checks}

Step 2 is split into three levels, and each level knows one thing.

\begin{description}[nosep,style=nextline]
  \item[A \texttt{Check} is a rule, and nothing else.] It is handed a list of readings and
        answers with the one reading at fault, which field to clear, and a sentence saying
        why --- or \texttt{None}. It never writes to the journal and never changes a reading.
        It is told at construction which field it judges and what to call it, so the same rule
        serves pH, temperature, weight or current without being rewritten.
  \item[An \texttt{Inspector} owns one part of a run.] It picks that part's readings out of the
        run, sorts them oldest first, builds the checks that part is held to, and does the
        clearing and the journalling for all of them.
  \item[The \texttt{Step} owns the loop over runs.] \texttt{SkipSuspectReadings} does nothing
        but put every run past every inspector.
\end{description}

The split is there so that a rule is written once and judged in one place. ``This value is
nothing like the rest of its own series'' is the same arithmetic for a tank weight and for the
supply current, and \texttt{ValueFarFromTheMedian} is that arithmetic, used twice with
different fields and different tolerances. It also keeps the thresholds out of the rules: no
check holds a number of its own, so a test can hand a check different numbers and watch it
behave differently without touching it.

There are four inspectors:

\begin{itemize}[nosep]
  \item \texttt{ConcTankInspector}
  \item \texttt{DiluteTankInspector}
  \item \texttt{RinseTankInspector}
  \item \texttt{PowerLogInspector}
\end{itemize}

\noindent The three tank ones share a parent,
\texttt{TankInspector}, which holds the construction of every check a tank may be given, so a
tank's own class is nothing but the list of what that tank is held to. As the code stands all
three tank lists are identical --- the same six checks in the same order. They are three
classes rather than one because the rules for one tank may well diverge from another's: the
dilute is topped up during a run and the concentrate is not, and the day a weight rule has to
differ between them, the change is a line in one class instead of an \texttt{if} inside a
shared one.

Each tank is also judged only against its own readings. A conductivity that is ordinary in the
dilute is impossible in the rinse, so there is no cross-tank comparison anywhere.

\subsection{What each check looks for}

Six check classes do all the judging. Three of them compare a reading with its own series,
two compare two instruments reading the same quantity, and one needs no comparison at all.

\subsubsection*{\texttt{ValueOutsideTheRange}}

Built for a value that is wrong on its own, whatever its neighbours say. It takes a low and a
high bound and rejects the first reading outside them:
\begin{equation}
  \label{eq:outside-range}
  \text{reject } x_i \quad\text{when}\quad \neg\,(\,\ell \le x_i \le h\,).
\end{equation}
It needs no minimum number of readings, because the judgement does not depend on any other
reading. It returns the \emph{earliest} offender rather than the biggest; order does not
matter here, because the inspector asks again after each one and works through them all. It
guards pH and temperature in every tank.

\subsubsection*{\texttt{ValueFarFromTheMedian}}

Built for a value that holds one level all run. With $m$ the median of the series and $s$ the
share allowed, it returns the reading furthest from $m$, if that distance is too far:
\begin{equation}
  \label{eq:far-from-median}
  i^\star = \arg\max_i\,\lvert x_i - m\rvert,
  \qquad
  \text{reject } x_{i^\star} \text{ when } \lvert x_{i^\star} - m\rvert > s\,\lvert m\rvert.
\end{equation}
The median is used rather than the mean because one bad value cannot drag it, and because it
does not care whether the bad value sits at the start of the run or the end. The tolerance is a
share of the median, not an absolute amount, so one number works for a rinse tank holding some
$4\unit{kg}$ and a concentrate tank holding some $2.1\unit{kg}$. It guards the tank weights and
the supply current, and it is the only check given the ``fewest readings'' threshold: with
fewer than three readings it says nothing, because with two each looks as odd as the other.

\subsubsection*{\texttt{ValueMovingTooFast}}

Built for a value that drifts but cannot lurch. Readings in this data are 5 to 240 minutes
apart, so the allowance is a rate and not a step --- strict on readings taken close together,
forgiving of slow drift across a long gap. For every consecutive pair it compares the speed of
the change with the allowance $a$ over $T$ minutes:
\begin{equation}
  \label{eq:moving-too-fast}
  \frac{\lvert x_{i+1} - x_i\rvert}{t_{i+1} - t_i} > \frac{a}{T}.
\end{equation}
Pairs no time apart are skipped, so the division is always safe. The fastest offending pair
wins. A pair says that two readings cannot both be right but not which one is wrong, so the
blame goes to whichever of the two sits further from the median of the series. It guards pH and
temperature in every tank, and needs at least two readings.

\subsubsection*{\texttt{JumpThatComesBack}}

Built for conductivity, which climbs or falls all run. There is no one level to compare
against, so the median is useless as a reference --- but a digit written down wrong still has a
signature: the series leaves its track and the next reading puts it straight back. A real
change of level stays changed. Over every window of three readings, with
$J = x_{i+1} - x_i$ and $B = x_{i+2} - x_{i+1}$,
\begin{equation}
  \label{eq:jump-and-back}
  \lvert J \rvert > g\,\lvert m \rvert
  \qquad\text{and}\qquad
  \frac{-B}{J} \ge c,
\end{equation}
and the middle reading $x_{i+1}$ is the one cleared. The first condition keeps an ordinary
wobble on a smooth series from counting as a jump; the median $m$ is used only to set that
scale. The second condition is the share of the jump that came back, and it is reported as a
percentage in the journal line. It can exceed 100\%: run \texttt{AR}'s concentrate at minute 60
put 132\% back, meaning the series overshot its way home. Fewer than three readings and the
check says nothing, since the window needs three. It is the one check with its own hard-coded
minimum rather than the shared threshold.

\subsubsection*{\texttt{PairedValuesTooFarApart}}

The stack voltage falls all run, so again there is no level to compare against --- but it was
written down twice, once from the supply and once from the multimeter. Those two agree to
hundredths of a volt, so a gap of a whole volt means one of them has a digit too many or too
few. The check takes the widest gap and rejects it if it is over the limit:
\begin{equation}
  \label{eq:voltage-gap}
  \text{reject the widest } \lvert v^{\text{supply}}_i - v^{\text{meter}}_i\rvert
  \quad\text{when it exceeds } \Delta_{\max}.
\end{equation}

\subsubsection*{\texttt{PairedValuesOutOfProportion}}

The same idea for the currents, which cannot be subtracted: the supply reads about a seventh
higher than the multimeter throughout, and in different units. So the two are held against
each other as a ratio, after the supply's amps are scaled by 1000 to meet the multimeter's
milliamps:
\begin{equation}
  \label{eq:current-ratio}
  r_i = \frac{1000\,I^{\text{supply}}_i}{I^{\text{meter}}_i}.
\end{equation}
The reading rejected is the one whose $r_i$ falls furthest from 1 among those outside
$[\,\ell,\,h\,]$. Readings where the multimeter wrote a zero are left out, so the ratio is
never divided by zero.

\subsubsection*{Which of two instruments gets the blame}

Both paired checks share \texttt{PairedValueCheck}, and both face the same problem: a
disagreement does not say which instrument is wrong. What does say is each column's own
history. The honest instrument stays near what it has been reading all run, so the blame goes
to whichever of the two sits further from the median of its \emph{own} column, measured as a
share of that median so that amps and milliamps compare fairly:
\begin{equation}
  \label{eq:odd-one}
  \text{blame} = \arg\max_{k \in \{\text{supply},\,\text{meter}\}}
  \frac{\lvert v^k_i - \operatorname{median}(v^k) \rvert}{\lvert \operatorname{median}(v^k)\rvert}.
\end{equation}
This is why the paired checks are the only ones that can clear a field other than the one they
are named after, and why their journal lines name the instrument. Run \texttt{AG} at minute 225
read $100.06\unit{V}$ on the supply against $10.00\unit{V}$ on the meter: a gap of
$90.06\unit{V}$, and the supply was the one nowhere near its own usual $11\unit{V}$, so the
supply's value was cleared. At minute 95 of the same run the supply read $11.20\unit{V}$ and the
meter $1.16\unit{V}$; this time the meter was the stranger in its own column, so the meter's
value went instead. Same check, same run, opposite verdicts, decided entirely by which column
the reading was out of step with.

For both paired checks a reading counts as judgeable only when \emph{both} instruments were
written down. One instrument alone is a reading with nothing to check it against, so it is
left alone.

\subsection{The numbers the checks judge by}

Every number a check compares against lives in one frozen dataclass called
\texttt{Thresholds}, with the reason for each one in the comment beside it. Change a number
there and every check that uses it changes with it. Table~\ref{tab:thresholds} is that file.

\begin{table}[htbp]
  \centering
  \caption{Every value in \texttt{thresholds.py}, and why it is what it is.}
  \label{tab:thresholds}
  \footnotesize
  \begin{tabular}{@{}llp{0.48\textwidth}@{}}
    \toprule
    Name & Value & Why \\
    \midrule
    \texttt{fewest\_readings\_to\_judge} & 3
      & Below this a check says nothing. With two readings each looks as odd as the other. \\
    \texttt{most\_of\_a\_series\_removable} & 0.34
      & The most of one series a check may clear. Needing more than a third means the file
        is wrong, not the readings. Past half, the median starts following the bad values. \\
    \texttt{impossible\_weight\_share} & 0.10
      & How far a tank weight may sit from its tank's usual weight. The three known bad
        weights are 90\% or more away; every good one is within 5\%. \\
    \texttt{most\_ph\_change} & 1.5
      & How far pH may move \ldots \\
    \texttt{most\_ph\_change\_in\_minutes} & 15.0
      & \ldots\ and over how long. A rate, because readings are 5 to 240 minutes apart. \\
    \texttt{plausible\_ph} & $(2.0,\,12.0)$
      & The range a pH reading must sit in whatever its neighbours say. These solutions cannot
        be at pH 0.07. \\
    \texttt{most\_temp\_change} & 2.0
      & How far temperature may move \ldots \\
    \texttt{most\_temp\_change\_in\_minutes} & 15.0
      & \ldots\ and over how long. The fastest real change in the data is $0.9\unit{^\circ C}$
        per quarter hour, so $2\unit{^\circ C}$ is a wide allowance. \\
    \texttt{plausible\_temp\_c} & $(5.0,\,40.0)$
      & The water sits at room temperature: every reading in the data is between 19.9 and
        $23.0\unit{^\circ C}$. \\
    \texttt{smallest\_jump\_share} & 0.10
      & How big a conductivity jump must be to count, as a share of what the series usually
        reads. Below a tenth, an ordinary wobble looks like a jump. \\
    \texttt{comes\_back\_share} & 0.7
      & How much of the jump the next reading must put back. At 0.7 this finds seven readings,
        every one a digit too many or too few. \\
    \texttt{impossible\_current\_share} & 0.10
      & How far the supply current may sit from its run's usual reading. It is held to
        within 2.2\% in nineteen of twenty-one runs. \\
    \texttt{largest\_voltage\_gap\_v} & 1.0
      & How far the two voltages may be apart. They agree to about $0.03\unit{V}$, so a gap
        this wide means a digit was written down wrong. \\
    \texttt{plausible\_current\_ratio} & $(0.62,\,1.6)$
      & The band the supply-to-meter current ratio must sit in. The supply reads about a
        seventh higher throughout, so a ratio and not a subtraction. \\
    \bottomrule
  \end{tabular}
\end{table}

Two of these are not any one check's. The fewest-readings number is handed only to the two
\texttt{ValueFarFromTheMedian} checks. The removable share belongs to the inspector, and no
check ever sees it.

\subsection{The one line where a value is cleared}

In the whole of the preprocessing there is exactly one place a reading's value is changed, in
\texttt{Inspector.\_clear\_what\_it\_rejects}:

\begin{quote}
\begin{verbatim}
reading, field_name, reason = worst
journal.skip(reading, field_name, reason)
setattr(reading, field_name, None)
\end{verbatim}
\end{quote}

That is the point of the design. The journal write is the line before the clearing, in the same
three lines, so there is no path through the code that clears a value without recording it.
Checks cannot clear anything --- they only answer with what they found --- so there is nowhere
else a value could quietly go missing. The journal entry takes its id and its record type from
the record itself rather than from the call site, so the id in the report and the reading that
was changed can never disagree.

\subsection{Asking a check again after every value it rejects}

A check is not run once over a series. It is asked for the single worst offender it can see;
if it finds one, that value is journalled and cleared, and the check is asked again from the
top. Two things make that the right loop.

The first is simply that a check answers with one reading, not a list, so a series holding two
bad values needs two passes. Run \texttt{AS}'s supply current held two: $0.46\unit{A}$ at
minute 210 and $0.44\unit{A}$ at minute 480, against a level of $0.045\unit{A}$. One pass would
have caught one of them.

The second is that the level a check compares against is recomputed each time, so it is a
better level once a bad value is gone. Every median, rate and ratio is worked out from the
values still present, which is what makes it safe to judge the second-worst reading against a
series the worst one has already left. How much it matters depends on the series. On
\texttt{AS}'s current the median read $0.045\unit{A}$ on both passes, because 57 readings do
not care about two of them. On run \texttt{AG}, where three voltage values went, the supply's
median moved only from $11.215$ to $11.19\unit{V}$ --- but the blame did move: the first pass
cleared the supply's $100.06\unit{V}$, and the two passes after it cleared the multimeter
instead.

The loop needs a stop. \texttt{Inspector.\_find\_clearing\_limit} sets how many times a single
check may fire on a single series:
\begin{equation}
  \label{eq:clearing-limit}
  \text{limit} = \max\!\left(1,\ \left\lfloor
    n \times \texttt{most\_of\_a\_series\_removable} \right\rfloor \right),
\end{equation}
where $n$ is the number of judgeable readings in the series \emph{before} any clearing. With
0.34, a pH series of three readings allows one removal, a typical tank series of seventeen
allows five, and a typical power log of forty-one rows allows thirteen. The $\max(1, \cdot)$
means a check can always fire at least once, even on the shortest series it will judge.

Two honest consequences. The first is the reason for the number: beyond about half a series the
median stops describing the good readings and starts describing the bad ones, at which point
the check turns around and clears the readings that were right. Stopping at a third keeps well
clear of that. If a third is not enough, the right conclusion is that the file is wrong, not
that the readings are.

The second is that \textbf{hitting the cap is silent}. The loop is a plain
\texttt{for \_ in range(limit)}; when it runs out of iterations it simply ends. Nothing is
journalled, nothing is printed, and the report will not tell you that a check still had
complaints when it was cut off. On the data as it stands this never happens: re-running the
step with the share raised to 0.5 and to 0.9 clears the same 27 values, so no check is being
held back. But that is a property of today's data, not a guarantee, and it is the one place in
this stage where something is left out without a line in the journal.

\subsection{Dropping a run against skipping a value}

The journal has two verbs, and they mean different things.

A \textbf{drop} is about a whole run. The run leaves the list and nothing downstream sees it.
A drop is written by \texttt{Preprocessor} itself, not by the step: after each step it compares
the run codes it handed in with the run codes it got back, and journals every run that did not
come back, with the step's own name as the reason. That is why a drop line reads
\emph{no lab results at both ends} rather than a sentence --- the step name \emph{is} the
reason, and the step does not have to remember to report anything.

A \textbf{skip} is about one value. The reading stays in the run with everything else it
carries, and one field is set to \texttt{None}.

Setting the field to \texttt{None} rather than deleting the reading matters, for three reasons.
A \texttt{TankReading} carries conductivity, pH, temperature and weight together, and may carry
the lab sample drawn at that minute --- throwing the row away over one bad pH would throw away
good values and possibly a lab result that cost money to get. \texttt{None} is already a legal
state for every one of those fields, because the loader has to represent a blank cell in the
CSV the same way, so nothing downstream needs new handling for a cleared value. And
\texttt{Check.find\_judgeable\_readings} skips \texttt{None} values by definition, so a cleared
value is invisible to every later check and to every later median without any bookkeeping: it
is simply a reading that was never taken.

\subsection{What it does to the real runs}

Running \texttt{uv run limodel analyze} on the project data gives 32 journal lines before the
results table: 27 skipped readings and 5 dropped runs. All 32 come from this stage --- the
calculating and the fitting share the same journal and wrote nothing to it on this data. The
loader hands over 21 runs carrying 1049 tank readings and 891 power-log rows, and the stage
returns 16 runs --- the line the report prints is \texttt{using 16 of 21 runs}.

Step 1 drops nothing, because all 21 runs are marked for use. All five drops come from step 3:
runs \texttt{AU}, \texttt{AV}, \texttt{AW}, \texttt{AX} and \texttt{AY} do not have both metals
reported at both ends of both tanks. Step 2 drops nothing by construction and accounts for all
27 skips --- 13 tank values and 14 power-log values, out of 1940 readings.

Table~\ref{tab:skips-by-check} is which check found them.

\begin{table}[htbp]
  \centering
  \caption{The 27 skipped values by the check that rejected them, and by the field cleared.}
  \label{tab:skips-by-check}
  \begin{tabular}{@{}lrlr@{}}
    \toprule
    Check & Values & Field cleared & Values \\
    \midrule
    \texttt{PairedValuesTooFarApart}     & 9 & \texttt{meter\_voltage\_v}   & 8 \\
    \texttt{JumpThatComesBack}           & 7 & \texttt{conductivity\_us\_cm} & 7 \\
    \texttt{ValueFarFromTheMedian}       & 6 & \texttt{ph}                  & 3 \\
    \texttt{PairedValuesOutOfProportion} & 2 & \texttt{weight\_g}           & 3 \\
    \texttt{ValueOutsideTheRange}        & 2 & \texttt{supply\_current\_a}  & 3 \\
    \texttt{ValueMovingTooFast}          & 1 & \texttt{meter\_current\_ma}  & 2 \\
                                         &   & \texttt{supply\_voltage\_v}  & 1 \\
    \midrule
    Total                                & 27 &                             & 27 \\
    \bottomrule
  \end{tabular}
\end{table}

The journal lines themselves say what happened in plain words. Three of the real ones:

\begin{quote}\footnotesize
\begin{verbatim}
skipped  reading  AJ:conc:0
    weight 209.7 g is 1802.15 g away from the usual 2011.85 g
skipped  reading  AS:dil:255
    conductivity jumped to 173997 and the next reading
    put 100% of it back
skipped  reading  AG:225
    the two voltage readings differ by 90.06 V, and the
    supply's 100.06 V is unlike its own column
\end{verbatim}
\end{quote}

Read them as three different shapes of mistake. The first is a decimal point in the wrong
place: $209.7\unit{g}$ in a tank whose next recorded weight is $2098.9\unit{g}$ and whose usual
weight is about $2012\unit{g}$. The second is a digit too many, and it is only findable because
the series carries straight on afterwards: $17549.7$, then $173997.1$, then $17307.2$. Nothing
here has to know what the right value was, only that the series left its track for exactly one
reading. The third is a supply voltage of $100.06\unit{V}$ where the stack was running near
$11\unit{V}$, and the multimeter stayed on its own track, so the supply column is the one that
slipped.

The id tells you where to look in the sheet. An id of \texttt{AJ:conc:0} means run
\texttt{AJ}, concentrate tank, minute 0. An id of \texttt{AG:225} has no tank in it because it
is a power-log row: run \texttt{AG}, minute 225. Every line is a cell to go and check, which is
the whole purpose of keeping the journal rather than cleaning silently.

\section{Working out what each run came to}

Preprocessing hands on a list of runs the lab can answer for. This stage turns each one into a
single row of numbers. The class that drives it is \texttt{Calculator}, in
\path{src/limodel/analyzer/calculate/calculator.py}, and its one public method is
\texttt{find\_results(runs)}. Every row it produces is a \texttt{Result}: a frozen dataclass
holding the run code, the membrane, the two settings the run was done at, the two outcomes the
model is fitted to, and the five further fields the report calls diagnostics. One of those five,
the highest voltage, is fitted as well, because a gate is placed on it.

Nothing in a \texttt{Result} is copied out of a CSV file. Every field is worked out from tank
weights, lab concentrations and the power log. The chain below is given in dependency order,
because that is the order the code reads in: the later quantities call the earlier ones rather
than re-deriving them.

Each quantity is one small class with static methods and no state. They are in
\path{src/limodel/analyzer/calculate/quantities/}, one file per class, and the file name is the
class name. The two constants they all lean on sit in \texttt{constants.py}:
$F = 96485\unit{C/mol}$, $M_\mathrm{Li} = 6.94\unit{g/mol}$,
$M_\mathrm{Mg} = 24.305\unit{g/mol}$, and a solution density of $1000\unit{g/L}$.

\subsection{Volume in a tank}

\textbf{\texttt{VolumeInATank}} --- \texttt{find\_litres(reading, tare\_g)} $\rightarrow$
litres.

Nothing on the rig measured volume. Each tank sat on a balance for the whole run and was weighed
empty beforehand, so the liquid is the balance reading less that empty weight:

\begin{equation}
  V(t) \;=\; \frac{w(t) - w_\mathrm{tare}}{\rho},
  \qquad \rho = 1000\unit{g/L}
  \label{eq:volume}
\end{equation}

$w(t)$ is \texttt{weight\_g} on the \texttt{TankReading} for that minute, and $w_\mathrm{tare}$
comes from the run's own \texttt{tare\_conc\_g}, \texttt{tare\_dil\_g} or \texttt{tare\_rinse\_g}
column, picked by \texttt{TARE\_FIELD\_BY\_TANK}. Taking $\rho$ as exactly
$1000\unit{g/L}$ is a deliberate simplification: these are dilute salt solutions, so a litre
weighs a kilogram closely enough that grammes on the balance can be read straight off as litres.
Returns \texttt{None} when the balance gave nothing for that minute, or the tare column is empty.

\subsection{Amount of a metal in a tank}

\textbf{\texttt{AmountInATank}} --- \texttt{find\_milligrams(reading, tare\_g, metal)}
$\rightarrow$ mg.

The lab reports a concentration, which on its own says nothing useful here. Every tank gains and
loses liquid through the run, so a concentration can rise while the metal in the tank falls. Only
the mass can be compared minute to minute:

\begin{equation}
  m_X(t) \;=\; c_X(t)\,V(t),
  \qquad [\unit{mg/L}] \times [\unit{L}] = [\unit{mg}]
  \label{eq:amount}
\end{equation}

$c_X(t)$ is read off the \texttt{Sample} hanging on that reading, through the \texttt{Metal}'s
own \texttt{field\_name} --- \texttt{li\_mg\_l} for lithium, \texttt{mg\_mg\_l} for magnesium.
Returns \texttt{None} when the minute has no sample at all, when the lab left that metal's cell
blank, or when \eqref{eq:volume} could not be worked out.

\subsection{Amount that crossed}

\textbf{\texttt{AmountThatCrossed}} --- \texttt{find\_milligrams(run, tank, metal)}
$\rightarrow$ mg.

How much of a metal a tank ended with, less how much it started with:

\begin{equation}
  \Delta m_X^{\,\mathrm{tank}} \;=\; m_X(t_\mathrm{last}) - m_X(t_\mathrm{first})
  \label{eq:crossed}
\end{equation}

Only sampled minutes can take part, since \eqref{eq:amount} needs a lab result. The ends are
found by \emph{minute}, not by position in the list: \texttt{find\_amounts\_by\_minute} builds a
dictionary keyed by minute, keeps only the minutes that produced a number, and the ends are taken
from \texttt{sorted()} of those keys. That matters because readings arrive in file order and
because preprocessing may have removed a minute in the middle --- or at an end, as it did on
AJ's concentrate at minute~0.

\texttt{find\_ends} returns the pair $(m_X(t_\mathrm{first}), m_X(t_\mathrm{last}))$ rather than
the difference, because recovery, the feed ratio and the balance gap all need the \emph{start}
on its own. Both methods return \texttt{None} when fewer than two sampled minutes of that tank
came through with a number, and \texttt{find\_ends} also returns \texttt{None} for a tank name
it does not know.

\subsection{Lithium recovery}

\textbf{\texttt{LithiumRecovery}} --- \texttt{find\_percent(run)} $\rightarrow$ per cent.

The share of the lithium the feed started with that ended up in the concentrate:

\begin{equation}
  R_\mathrm{Li} \;=\; 100 \times
  \frac{\Delta m_\mathrm{Li}^{\,\mathrm{conc}}}{m_\mathrm{Li}^{\,\mathrm{dil}}(t_\mathrm{first})}
  \label{eq:recovery}
\end{equation}

Milligrams stand on both the top and the bottom, so $M_\mathrm{Li}$ cancels. There is nothing to
gain by converting to moles, and one fewer place for a unit slip to hide. The result is a per
cent, not a fraction, because that is how the gate is written and how the report prints it.

The bottom is the dilute's \emph{first} sampled minute, not what it ended with. What the feed
held before anything moved is the whole prize; measuring against the end would make a run that
moved nothing look like a run that moved everything.

The top is a difference, \eqref{eq:crossed}, and that subtraction is not cosmetic. The
concentrate never starts with zero lithium, and across these runs its starting amount spans a
factor of more than five hundred, depending on which salt the concentrate was made up with.
Over the sixteen runs the calculator uses it runs from $2.6\unit{mg}$ up to $1455.1\unit{mg}$.
The low end is AN and AO, whose concentrate was made up on a sodium salt rather than a lithium
one, so only traces were there at the start ($18.6$ and $2.6\unit{mg}$). The high end is AP and
AQ, started on LiCl at $5\unit{g/L}$ ($1414.4$ and $1455.1\unit{mg}$). The twelve runs started
on LiCl at $0.5\unit{g/L}$ sit between $126.4$ and $161.7\unit{mg}$. On AQ the lithium already
in the
concentrate is about double what the whole feed held, so treating the start as zero would not
merely bias recovery --- it would report roughly $214\,\%$ in place of $11.4\,\%$.

As a check on \eqref{eq:recovery}, run AE: the concentrate gained
$309.55\unit{mg}$ of lithium, the dilute began with $808.74\unit{mg}$, and the report prints
$38.3\,\%$.

\subsection{Charge that passed}

\textbf{\texttt{ChargeThatPassed}} --- \texttt{find\_coulombs(run)} $\rightarrow$ C, and
\texttt{find\_mah\_per\_cm2(run)} $\rightarrow$ mA\,h/cm\textsuperscript{2}.

No coulomb counter was fitted. The supply current was written down every so many minutes, so the
charge is the area under those readings, by the trapezium rule over consecutive logged minutes:

\begin{equation}
  Q \;=\; 60 \sum_{k} \frac{I_k + I_{k+1}}{2}\,\bigl(t_{k+1} - t_k\bigr)
  \label{eq:charge}
\end{equation}

The $60$ converts amp-minutes to amp-seconds, and an amp-second is a coulomb. Spread over the
membrane area and converted to the unit the field quotes:

\begin{equation}
  q \;=\; \frac{Q}{3600\,A} \times 1000
  \qquad [\unit{mA\,h/cm^2}]
  \label{eq:charge-density}
\end{equation}

The gaps $t_{k+1}-t_k$ are whatever the log holds; nothing is resampled onto an even grid, and no
reading is carried forward past the last logged minute. Only minutes with a
\texttt{supply\_current\_a} take part, so a minute preprocessing removed simply widens the
trapezium either side of it. \texttt{find\_coulombs} returns \texttt{None} on fewer than two such
minutes, and \texttt{find\_mah\_per\_cm2} also returns \texttt{None} when
\texttt{area\_cm2} is zero or missing. AE gives $Q = 1178.55\unit{C}$ over
$A = 6\unit{cm^2}$, which is $54.6\unit{mA\,h/cm^2}$.

\subsection{Current density}

\textbf{\texttt{CurrentDensity}} --- \texttt{find\_ma\_per\_cm2(run)} $\rightarrow$
mA/cm\textsuperscript{2}.

\begin{equation}
  J \;=\; \frac{1000 \times \operatorname{median}\bigl(\{I_k\}\bigr)}{A}
  \qquad [\unit{mA/cm^2}]
  \label{eq:current-density}
\end{equation}

This is the first of the two columns the polynomial is fitted over, so how it is obtained
matters more than it looks. It comes from the \emph{measured} supply current, not from the set
point: the \texttt{set\_current\_ma} column is empty on all twenty-one runs in \path{data/runs.csv},
so there is no set point to read. The median rather than the mean, because a run was held at one
current throughout and the number that labels the whole run should not be dragged by a single
misread minute. Returns \texttt{None} when no minute carries a current, or the run has no area.

\subsection{Stack voltage}

\textbf{\texttt{StackVoltage}} --- \texttt{find\_highest\_volts(run)} $\rightarrow$ V, and
\texttt{find\_starting\_volts(run)} $\rightarrow$ V.

\begin{equation}
  U_\mathrm{max} \;=\; \max_k U_k,
  \qquad
  U_\mathrm{start} \;=\; U_{k_\mathrm{min}}
  \label{eq:voltage}
\end{equation}

$U_\mathrm{max}$ is the highest \texttt{supply\_voltage\_v} the run reached, and it is the figure
the $20\unit{V}$ gate is applied to: a ceiling is only cleared if it is cleared at the worst
minute, not on average. $U_\mathrm{start}$ is the voltage at the earliest minute that has one ---
the earliest, deliberately, not the lowest --- and says what the stack asked for before anything
had crossed. Only $U_\mathrm{max}$ reaches a \texttt{Result}; $U_\mathrm{start}$ is there for
reading a run. Both return \texttt{None} when the power log never recorded a voltage.

\subsection{Transport numbers}

\textbf{\texttt{TransportNumbers}} --- \texttt{find\_numbers(run)} $\rightarrow$
$(t_\mathrm{Li}, t_\mathrm{Mg})$, dimensionless, summing to $1$.

A transport number is a share of the \emph{charge} carried, not a share of the ions. So each
metal's crossing is turned into millimoles and then weighted by the charge its own ion carries:
$z_\mathrm{Li} = 1$, $z_\mathrm{Mg} = 2$. Those charges live on the \texttt{Metal} dataclass
alongside the molar mass, which is why they travel together.

\begin{equation}
  n^\star_X \;=\; \frac{z_X\,\Delta m_X^{\,\mathrm{conc}}}{M_X},
  \qquad
  t_\mathrm{Li} = \frac{n^\star_\mathrm{Li}}{n^\star_\mathrm{Li} + n^\star_\mathrm{Mg}},
  \qquad
  t_\mathrm{Mg} = \frac{n^\star_\mathrm{Mg}}{n^\star_\mathrm{Li} + n^\star_\mathrm{Mg}}
  \label{eq:transport}
\end{equation}

Why magnesium counts twice: the current through the stack is carried by charge, and one
Mg\textsuperscript{2+} crossing the membrane moves two units of charge while one
Li\textsuperscript{+} moves one. If the $z_X$ were dropped, these would be shares of the ion
count, and a membrane that let through equal moles of each would look even-handed when in fact
magnesium was doing two-thirds of the work. Returns \texttt{None} when either metal's crossing is
unknown, or when the two weighted amounts sum to zero. On AE,
$(t_\mathrm{Li}, t_\mathrm{Mg}) = (0.680, 0.320)$.

\subsection{The two feed ratios}

\textbf{\texttt{FeedRatio}} has two methods, both returning mol/mol:
\texttt{find\_li\_over\_mg\_in\_the\_dilute(run)} and
\texttt{find\_mg\_over\_li\_from\_the\_setup(run)}.

There are two, kept apart by name because they answer different questions and are not
reciprocals of one another.

The \emph{measured} one is taken from the lab, at the first dilute minute for which both metals
were reported:

\begin{equation}
  r_\mathrm{dil} \;=\;
  \frac{m_\mathrm{Li}^{\,\mathrm{dil}}(t_0)\,/\,M_\mathrm{Li}}
       {m_\mathrm{Mg}^{\,\mathrm{dil}}(t_0)\,/\,M_\mathrm{Mg}}
  \label{eq:feed-measured}
\end{equation}

$t_0$ is the earliest minute present in both metals' dictionaries --- the intersection of the
two, so a minute where the lab reported only one metal is skipped rather than half-used. Both
amounts come from the same tank at the same minute, so the litres in \eqref{eq:amount} cancel and
the milligrams carry the ratio on their own. This is the one permselectivity divides by.

The \emph{set up} one is taken from the run's own feed columns, and is the other way up:

\begin{equation}
  r_\mathrm{setup} \;=\;
  \frac{c_\mathrm{Mg}^{\,\mathrm{feed}}\,/\,M_\mathrm{Mg}}
       {c_\mathrm{Li}^{\,\mathrm{feed}}\,/\,M_\mathrm{Li}}
  \label{eq:feed-setup}
\end{equation}

using \texttt{mg\_feed\_mg\_l} and \texttt{li\_feed\_mg\_l}. This is the second of the two
columns the polynomial is fitted over, and it is the intended ratio rather than the measured one
on purpose: it is a dial setting that was chosen before the run, known exactly, and the same for
every replicate at that point. Using a measured value would put lab noise into a design column
and spread what should be one design point into a cloud. It goes into \texttt{Result} as
\texttt{feed\_mg\_li\_ratio}, which is why the report's column is headed Mg/Li. Each returns
\texttt{None} when its inputs are missing or its denominator is zero. On AE,
$r_\mathrm{dil} = 0.852$ and $r_\mathrm{setup} = 1.428$.

\subsection{Permselectivity}

\textbf{\texttt{Permselectivity}} --- \texttt{find\_value(run)} $\rightarrow$
dimensionless.

The outcome the whole project maximises:

\begin{equation}
  P \;=\; \frac{t_\mathrm{Li} / t_\mathrm{Mg}}{r_\mathrm{dil}}
  \label{eq:perm}
\end{equation}

The top is what the membrane delivered. The bottom is what it was offered. Dividing by
$r_\mathrm{dil}$ is what makes the number mean anything across runs, and the arithmetic is worth
following. Writing $n_X$ for the moles of each metal that crossed, \eqref{eq:transport} gives

\begin{equation}
  \frac{t_\mathrm{Li}}{t_\mathrm{Mg}} \;=\;
  \frac{z_\mathrm{Li}\,n_\mathrm{Li}}{z_\mathrm{Mg}\,n_\mathrm{Mg}} \;=\;
  \frac{n_\mathrm{Li}}{2\,n_\mathrm{Mg}}
  \label{eq:transport-ratio}
\end{equation}

Now take a membrane with no preference at all, letting the two metals through in proportion to
what the feed holds: $n_\mathrm{Li}/n_\mathrm{Mg} = r_\mathrm{dil}$. Then
\eqref{eq:transport-ratio} is $r_\mathrm{dil}/2$, and \eqref{eq:perm} comes to exactly $0.5$ ---
whatever the feed was. That is the number to read $P$ against: above $0.5$ the membrane preferred
lithium, below it magnesium, and the factor of two is the charge weighting, not a choice.

Without the division the baseline would slide. The same indifferent membrane would score
$r_\mathrm{dil}/2$: that is $0.426$ at AE's feed, and $0.043$ at a feed ten times richer in
magnesium. Two runs at different feeds could not then be set side by side. Because the feed ratio
is one of the two model inputs, this is not a detail. The polynomial is asked how $P$ changes as
the feed ratio moves; if the baseline moved with it too, the fitted feed-ratio term would be
measuring the normalisation rather than the membrane, and every conclusion about feed sensitivity
would be an artefact.

Returns \texttt{None} when the transport numbers or the measured feed ratio are missing, or when
$t_\mathrm{Mg}$ or $r_\mathrm{dil}$ is zero. On AE,
$P = (0.680/0.320)/0.852 = 2.498$, about five times indifference.

\subsection{Energy per mole of lithium}

\textbf{\texttt{EnergyPerMoleOfLithium}} --- \texttt{find\_joules(run)} $\rightarrow$ J, and
\texttt{find\_watt\_hours(run)} $\rightarrow$ W\,h/mol.

Power is multiplied minute by minute before any averaging, because voltage and current do not
move together and the average of a product is not the product of the averages. The same
trapezium rule as \eqref{eq:charge} then runs over the minutes that carry both:

\begin{equation}
  E \;=\; 60 \sum_{k} \frac{U_k I_k + U_{k+1} I_{k+1}}{2}\,\bigl(t_{k+1} - t_k\bigr)
  \qquad [\unit{J}]
  \label{eq:energy}
\end{equation}

\begin{equation}
  n_\mathrm{Li} \;=\; \frac{\Delta m_\mathrm{Li}^{\,\mathrm{conc}}}{1000\,M_\mathrm{Li}},
  \qquad
  e \;=\; \frac{E}{3600\,n_\mathrm{Li}}
  \qquad [\unit{W\,h/mol}]
  \label{eq:energy-per-mole}
\end{equation}

Per mole of lithium delivered, rather than per hour or per litre, so a short run and a long one,
a gentle current and a hard one, can be compared on the thing being paid for. Note the set of
minutes here is narrower than in \eqref{eq:charge}: a minute with a current but no voltage counts
towards charge and not towards energy. \texttt{find\_joules} returns \texttt{None} on fewer than
two minutes carrying both. \texttt{find\_watt\_hours} also returns \texttt{None} when
$n_\mathrm{Li} \le 0$, since a concentrate that gained no lithium, or lost some, has no per-mole
cost to quote. AE comes to $E = 14759.8\unit{J}$ and $e = 91.9\unit{W\,h/mol}$.

\subsection{The balance gap}

\textbf{\texttt{BalanceGap}} --- \texttt{find\_percent(run, metal)} $\rightarrow$ per cent,
computed once for lithium and once for magnesium.

\begin{equation}
  L_X \;=\; m_X^{\,\mathrm{dil}}(t_\mathrm{first}) - m_X^{\,\mathrm{dil}}(t_\mathrm{last}),
  \qquad
  G_X \;=\; 100 \times
  \frac{L_X - \bigl(\Delta m_X^{\,\mathrm{conc}} + \Delta m_X^{\,\mathrm{rinse}}\bigr)}{L_X}
  \label{eq:gap}
\end{equation}

This corrects nothing. No other quantity reads it, no fit uses it, and no run is dropped on it.
It exists so a reader can judge how much weight one run's numbers deserve.

A large positive gap means the dilute lost metal that was not found in either of the other two
tanks --- the readings do not add up, and the run's recovery and permselectivity should be read
with suspicion. A negative gap means more was found than the dilute is recorded as having lost,
which usually points at a weight or a concentration being wrong at one end. AT shows
$-272.5\,\%$: its dilute barely changed, losing $5.05\unit{mg}$, while its concentrate gained
$18.82\unit{mg}$, so a small absolute mismatch divides by a tiny denominator. AT's recovery of
$2.3\,\%$ should be read in that light.

The liquid carried off in the sample bottles is \emph{not} accounted for, by choice. Every sample
drawn removes both volume and metal from a tank for good. Rather than model that removal --- which
would mean tracking bottle volumes nobody recorded reliably --- it is left out, and the gap
absorbs it. So a modest positive gap is expected, and the gap is a bound on the error rather than
a measurement of it. A missing rinse figure is likewise counted as zero gained, which can only
make the gap look worse; for a diagnostic, erring towards alarm is the honest direction. Fifteen
of the sixteen used runs do have a rinse figure. Returns \texttt{None} when the dilute ends or
the concentrate's gain cannot be worked out, or when $L_X \le 0$: a share of a loss means nothing
where there was no loss, and a dilute that gained would turn the sign of the gap over.

\subsection{Missing is never a guess}

Every method above returns \texttt{None} when it cannot be worked out. Nothing is filled in with
a zero, a mean, or a carried-forward value. A zero is a claim about the chemistry and would be
fitted as one; \texttt{None} is a claim about the records and can be counted.

\texttt{Calculator.\_find\_result} then applies one rule. Permselectivity and lithium recovery
are what the model is fitted to, so a run missing either has nothing to contribute: the drop is
written in the journal with the reason naming which of the two failed, and \texttt{None} comes
back instead of a \texttt{Result}. Everything else may be \texttt{None} on a run that is
otherwise sound, and that is not a reason to let the run go --- the report prints a dash in the
cell and names the run under \emph{diagnostics missing} at the end. On the current data all
sixteen surviving runs produce a complete row, so no drop happens at this stage and no dash is
printed; the five runs the report says were dropped were dropped during preprocessing.

\subsection{Every quantity in one table}

\begin{table}[htbp]
  \centering
  \caption{The quantities, in dependency order. Only the fields of \texttt{Result} reach the
    model or the report.}
  \label{tab:quantities}
  \small
  \begin{tabular}{@{}l l l p{0.30\linewidth}@{}}
    \toprule
    Quantity & Class & Unit & Used for \\
    \midrule
    Volume in a tank & \texttt{VolumeInATank} & $\unit{L}$
      & turning a concentration into a mass \\
    Amount in a tank & \texttt{AmountInATank} & $\unit{mg}$
      & one tank at one minute \\
    Amount that crossed & \texttt{AmountThatCrossed} & $\unit{mg}$
      & recovery, transport numbers, feed ratio, energy, balance gap \\
    Lithium recovery & \texttt{LithiumRecovery} & \%
      & fitted outcome; the recovery gate \\
    Charge that passed & \texttt{ChargeThatPassed} & $\unit{mA\,h/cm^2}$
      & diagnostic: how hard the run was driven in total \\
    Current density & \texttt{CurrentDensity} & $\unit{mA/cm^2}$
      & model input $J$; the setting the search optimises \\
    Stack voltage & \texttt{StackVoltage} & $\unit{V}$
      & fitted outcome; the $20\unit{V}$ gate \\
    Transport numbers & \texttt{TransportNumbers} & ---
      & permselectivity \\
    Feed ratio, measured & \texttt{FeedRatio} & $\unit{mol/mol}$
      & the denominator of permselectivity \\
    Feed ratio, set up & \texttt{FeedRatio} & $\unit{mol/mol}$
      & model input $R$; the second search setting \\
    Permselectivity & \texttt{Permselectivity} & ---
      & fitted outcome; what the search maximises \\
    Energy per mole Li & \texttt{EnergyPerMoleOfLithium} & $\unit{W\,h/mol}$
      & diagnostic: the electricity cost of a run \\
    Balance gap, Li and Mg & \texttt{BalanceGap} & \%
      & diagnostic only; corrects nothing \\
    \bottomrule
  \end{tabular}
\end{table}

\section{One run, start to finish}

Every formula in this guide is small. What is easy to lose is the chain: a balance reading
becomes litres, litres and a lab concentration become milligrams, two sets of milligrams
become a crossing, and only then does anything look like a result. This section carries run
AE through that chain with the numbers the code actually produces, so any step can be checked
on paper. Run AG follows in a table beside it.

Every figure below came from running the code, not from reading the formulas. A reader who
wants to repeat it can print any intermediate value the same way:

\begin{quote}
\texttt{uv run python -c "... AmountThatCrossed} \\
\texttt{.find\_milligrams(run, CONCENTRATE, LITHIUM)"}
\end{quote}

\subsection{What AE arrived with}

AE used a GO-10 cation membrane over $6.0\unit{cm^2}$, fed at $400\unit{mg/L}$ lithium and
$2000\unit{mg/L}$ magnesium, and ran for 240 minutes. Only the readings in
Table~\ref{tab:ae-inputs} reach the two outcomes the model is fitted to. Everything else in
AE's eighty-odd rows is there for the checks, for the plots, or for nothing at all.

\begin{table}[htbp]
\centering
\caption{The readings AE's result is built from. The tares come from the run sheet; the
weights from the balance under each tank; the concentrations from the lab.}
\label{tab:ae-inputs}
\begin{tabular}{llrr}
\toprule
 & & minute 0 & minute 240 \\
\midrule
Concentrate & tare & \multicolumn{2}{c}{$96.2\unit{g}$} \\
            & weight            & $2119.7\unit{g}$ & $2101.5\unit{g}$ \\
            & lithium           & $63.56\unit{mg/L}$ & $218.5\unit{mg/L}$ \\
            & magnesium         & $2.974\unit{mg/L}$ & $130.0\unit{mg/L}$ \\
\addlinespace
Dilute      & tare & \multicolumn{2}{c}{$111.5\unit{g}$} \\
            & weight            & $2117.8\unit{g}$ & --- \\
            & lithium           & $403.1\unit{mg/L}$ & --- \\
            & magnesium         & $1657\unit{mg/L}$ & --- \\
\bottomrule
\end{tabular}
\end{table}

The dilute's last minute is left out of the table on purpose. Recovery divides by what the
feed held \emph{before} anything moved, so only its first sampled minute is needed; its last
minute is used by the balance gap, which is a diagnostic. The two tares matter as much as the
weights do, because neither tank weight means anything until its own empty weight is taken
off it.

\subsection{Weights become litres}

Nothing measured a volume. Each tank sat on a balance all run, so
\texttt{VolumeInATank}'s \texttt{find\_litres} turns a reading into litres by subtracting
that tank's tare and dividing by $1000\unit{g/L}$ --- the dilute salt solutions in these tanks weigh
close enough to a kilogram per litre that the balance can be read as a volume.

\begin{align}
V &= \frac{w - w_{\mathrm{tare}}}{1000} \label{eq:volume-ae} \\
V_{\mathrm{conc}}(0)   &= \frac{2119.7 - 96.2}{1000}  = 2.0235\unit{L} \nonumber \\
V_{\mathrm{conc}}(240) &= \frac{2101.5 - 96.2}{1000}  = 2.0053\unit{L} \nonumber \\
V_{\mathrm{dil}}(0)    &= \frac{2117.8 - 111.5}{1000} = 2.0063\unit{L} \nonumber
\end{align}

The concentrate lost $18.2\unit{g}$ over the run, about nine parts in a thousand of its
volume. That is why \eqref{eq:volume-ae} is run at each sampled minute rather than once: a
concentration compared minute to minute against a volume that has shifted is comparing two
different amounts of liquid.

\subsection{Litres and concentrations become milligrams}

\texttt{AmountInATank.find\_milligrams} multiplies the lab's $\unit{mg/L}$ by the litres from
\eqref{eq:volume-ae}. Four numbers come out of AE's concentrate:

\begin{align}
m &= c \, V \label{eq:amount-ae} \\
\mathrm{Li}(0)   &= 63.56 \times 2.0235 = 128.614\unit{mg} \nonumber \\
\mathrm{Li}(240) &= 218.5 \times 2.0053 = 438.158\unit{mg} \nonumber \\
\mathrm{Mg}(0)   &= 2.974 \times 2.0235 =   6.018\unit{mg} \nonumber \\
\mathrm{Mg}(240) &= 130.0 \times 2.0053 = 260.689\unit{mg} \nonumber
\end{align}

The feed's first minute goes the same way: $403.1 \times 2.0063 = 808.740\unit{mg}$ of
lithium and $1657 \times 2.0063 = 3324.439\unit{mg}$ of magnesium.

\subsection{What crossed}

\texttt{AmountThatCrossed} takes the sampled minutes of one tank, sorts them \emph{by
minute}, and subtracts the first from the last. It sorts because the readings arrive in file
order, which is not guaranteed to be time order; taking the first and last rows as they came
would silently reverse the sign on a run whose rows were entered out of sequence.

\begin{align}
\Delta m_{\mathrm{Li}} &= 438.158 - 128.614 = 309.544\unit{mg}
   &&\rightarrow\; \frac{309.544}{6.94} = 44.603\unit{mmol} \label{eq:crossed-li} \\
\Delta m_{\mathrm{Mg}} &= 260.689 -   6.018 = 254.671\unit{mg}
   &&\rightarrow\; \frac{254.671}{24.305} = 10.478\unit{mmol} \label{eq:crossed-mg}
\end{align}

Read in milligrams, magnesium looks almost as mobile as lithium. Read in millimoles it does
not: lithium crossed in four times the molar quantity. That is the whole reason the code
converts before it compares, and it is why the next two subsections look so different from
each other.

\subsection{Recovery}

\texttt{LithiumRecovery.find\_percent} divides what reached the concentrate by what the feed
started with. Both sides are milligrams, so the molar mass cancels and there is nothing to
gain from going through moles.

\begin{equation}\label{eq:recovery-ae}
R_{\mathrm{Li}} = 100 \times \frac{\Delta m_{\mathrm{Li}}}{m_{\mathrm{dil,Li}}(0)}
                = 100 \times \frac{309.544}{808.740} = 38.275\,\%
\end{equation}

The denominator is the feed's \emph{first} minute, never its last. Against the
$505.257\unit{mg}$ the feed ended with, the same crossing would read $61.3\,\%$, which would
reward a run for losing feed to the sample bottles.

\subsection{The charge that passed}

No meter totalled the charge. The supply current was written down every five to fifteen
minutes, so \texttt{ChargeThatPassed} integrates those readings with the trapezoid rule:
average the two currents across a gap and multiply by its length.

\begin{equation}\label{eq:charge-ae}
Q = 60 \sum_{k} \frac{I_k + I_{k+1}}{2}\,(t_{k+1} - t_k)
\end{equation}

The factor of 60 is the only unit conversion in \eqref{eq:charge-ae}: the log counts minutes and
a coulomb is an amp-second. AE's log holds 30 minutes with a current, from minute 0 to minute
240. The first gap contributes $\tfrac{1}{2}(0.081 + 0.082) \times 10 = 0.815\unit{A\,min}$,
the last $0.082 \times 15 = 1.23\unit{A\,min}$, and all 29 gaps together come to
$19.6425\unit{A\,min}$. So

\begin{align}
Q &= 19.6425 \times 60 = 1178.55\unit{C} \label{eq:ae-charge} \\
\frac{Q}{A} &= \frac{1178.55}{3600 \times 6.0} \times 1000 = 54.563\unit{mA\,h/cm^2}
   \nonumber
\end{align}

Minute 170 of AE's log has a voltage but no current, so it never enters \eqref{eq:charge-ae} and
the curve bridges minute 165 to minute 180 in one step. The run's current density is a
separate number: \texttt{CurrentDensity} takes the \emph{median} of the same currents,
$0.082\unit{A}$, giving $0.082 \times 1000 / 6.0 = 13.667\unit{mA/cm^2}$. The median rather
than the mean, because one run holds one current all the way through and a single misread
minute should not move the label on the whole run.

\subsection{The transport numbers}

A transport number is a share of the \emph{charge}, not of the ions. Every magnesium ion
carries two charges, so \texttt{TransportNumbers} weights each metal's millimoles by the
charge on its own ion. The charge sits on the \texttt{Metal} object next to the molar mass,
which is what stops the weighting from being forgotten at a call site.

\begin{align}
z_{\mathrm{Li}} n_{\mathrm{Li}} &= 1 \times \frac{309.544}{6.94}   = 44.603 \label{eq:tn} \\
z_{\mathrm{Mg}} n_{\mathrm{Mg}} &= 2 \times \frac{254.671}{24.305} = 20.956 \nonumber \\
t_{\mathrm{Li}} &= \frac{44.603}{44.603 + 20.956} = \frac{44.603}{65.559} = 0.6803
   \nonumber \\
t_{\mathrm{Mg}} &= \frac{20.956}{65.559} = 0.3197 \nonumber
\end{align}

Lithium crossed in $4.26$ times the molar quantity magnesium did, yet carried only $2.13$
times the charge. Doubling magnesium in \eqref{eq:tn} halves that advantage, and the two
shares add to one by construction.

\subsection{The feed ratio and permselectivity}

$t_{\mathrm{Li}}/t_{\mathrm{Mg}}$ on its own says as much about the feed as about the
membrane: a feed with little magnesium in it cannot send much magnesium across. So
\texttt{Permselectivity} divides the transport ratio by the ratio the feed actually offered,
which \texttt{FeedRatio}'s \texttt{find\_li\_over\_mg\_\allowbreak in\_the\_dilute} measures
from the dilute's first sampled minute. Both amounts come from one tank at one minute, so
the litres cancel and the milligrams carry the ratio.

\begin{align}
r_{\mathrm{feed}} &= \frac{808.740 / 6.94}{3324.439 / 24.305}
                   = \frac{116.533}{136.780} = 0.85197 \label{eq:feed-ratio} \\
P &= \frac{t_{\mathrm{Li}} / t_{\mathrm{Mg}}}{r_{\mathrm{feed}}}
   = \frac{0.6803 / 0.3197}{0.85197} = \frac{2.12838}{0.85197} = 2.4982 \label{eq:perm-ae}
\end{align}

Dividing by the feed ratio pins ``no preference'' at $0.5$ whatever the feed was, which is
the only thing that lets AE and AG be set beside each other at all. AE's $2.4982$ says the
stack favoured lithium about five times over what the feed handed it.

Note that \eqref{eq:feed-ratio} is lithium over magnesium, measured in the tank. The report's
\texttt{Mg/Li} column is a different number. \texttt{FeedRatio}'s other method,
\texttt{find\_mg\_over\_li\_\allowbreak from\_the\_setup}, reads the run sheet's own feed
concentrations, the other way up, giving
$(2000/24.305)/(400/6.94) = 1.4277$ for AE. The two are kept apart by name because they
answer different questions --- the measured one is what \eqref{eq:perm-ae} divides by, the set-up
one is a design column of the polynomial.

\subsection{The voltage and the energy}

AE's power log never passed $16.0\unit{V}$, reached at minute 0, and fell to
$11.31\unit{V}$ by minute 240 as the concentrate gained salt and its resistance dropped. The
highest voltage is what the $20\unit{V}$ supply ceiling is judged against, and on a healthy
run it is the starting voltage.

Energy uses the same trapezoid as \eqref{eq:charge-ae}, over the product of volts and amps, and
only over minutes that carry \emph{both}. AE has 30 such minutes, totalling
$245.996\unit{W\,min}$:

\begin{align}
E &= 245.996 \times 60 = 14759.78\unit{J} = 4.0999\unit{W\,h} \label{eq:energy-ae} \\
n_{\mathrm{Li}} &= \frac{309.544}{1000 \times 6.94} = 0.044603\unit{mol} \nonumber \\
\frac{E}{n_{\mathrm{Li}}} &= \frac{4.0999}{0.044603} = 91.92\unit{W\,h/mol} \nonumber
\end{align}

Per mole of lithium delivered, rather than per hour or per litre, so a short run and a long
one, or a gentle current and a hard one, can be compared on the same axis. The denominator is
the concentrate's gain from \eqref{eq:crossed-li}, and a run whose concentrate gained no
lithium gets \texttt{None} here instead of a division by zero.

\subsection{Run AG beside it}

AG is AE with one setting changed. Same GO-10 membrane, same $6.0\unit{cm^2}$, same
$400\unit{mg/L}$ lithium feed, same $0.082\unit{A}$ held for 240 minutes --- and
$8000\unit{mg/L}$ magnesium instead of $2000$, which is four times the magnesium on a molar
basis too, $5.7108$ against $1.4277$.

\begin{table}[htbp]
\centering
\caption{AE and AG, every stage. Both ran GO-10 at $13.667\unit{mA/cm^2}$ for 240 minutes;
only the feed magnesium differs.}
\label{tab:ae-ag}
\begin{tabular}{lrr}
\toprule
 & AE & AG \\
\midrule
Feed magnesium ($\unit{mg/L}$)          & $2000$ & $8000$ \\
Feed Mg/Li from the run sheet           & $1.4277$ & $5.7108$ \\
\addlinespace
Concentrate tare ($\unit{g}$)           & $96.2$ & $120.6$ \\
Dilute tare ($\unit{g}$)                & $111.5$ & $128.8$ \\
Concentrate volume, minute 0 ($\unit{L}$)   & $2.0235$ & $2.0167$ \\
Concentrate volume, minute 240 ($\unit{L}$) & $2.0053$ & $1.9847$ \\
Dilute volume, minute 0 ($\unit{L}$)        & $2.0063$ & $2.0312$ \\
\addlinespace
Concentrate Li, minute 0 ($\unit{mg}$)   & $128.614$ & $139.999$ \\
Concentrate Li, minute 240 ($\unit{mg}$) & $438.158$ & $193.488$ \\
Concentrate Mg, minute 0 ($\unit{mg}$)   & $6.018$ & $0.190$ \\
Concentrate Mg, minute 240 ($\unit{mg}$) & $260.689$ & $615.455$ \\
Dilute Li, minute 0 ($\unit{mg}$)        & $808.740$ & $868.135$ \\
Dilute Mg, minute 0 ($\unit{mg}$)        & $3324.439$ & $14129.027$ \\
\addlinespace
Li crossed ($\unit{mg}$ / $\unit{mmol}$) & $309.544$ / $44.603$ & $53.489$ / $7.707$ \\
Mg crossed ($\unit{mg}$ / $\unit{mmol}$) & $254.671$ / $10.478$ & $615.265$ / $25.314$ \\
\addlinespace
Lithium recovery                         & $38.275\,\%$ & $6.161\,\%$ \\
Charge passed ($\unit{C}$)               & $1178.55$ & $1179.90$ \\
Charge passed ($\unit{mA\,h/cm^2}$)      & $54.563$ & $54.625$ \\
$t_{\mathrm{Li}}$ / $t_{\mathrm{Mg}}$    & $0.6803$ / $0.3197$ & $0.1321$ / $0.8679$ \\
Feed Li/Mg in the dilute                 & $0.85197$ & $0.21518$ \\
Permselectivity                          & $2.4982$ & $0.70745$ \\
\addlinespace
Highest voltage ($\unit{V}$)             & $16.0$ & $13.9$ \\
Energy ($\unit{W\,h/mol}$ of Li)         & $91.92$ & $474.68$ \\
Li balance gap                           & $-19.645\,\%$ & $18.406\,\%$ \\
Mg balance gap                           & $35.002\,\%$ & $44.626\,\%$ \\
\bottomrule
\end{tabular}
\end{table}

The two runs pushed essentially the same charge --- $1178.55$ against
$1179.90\unit{C}$, a tenth of a percent apart --- and spent it on completely different work.
Lithium recovery fell from $38.3\,\%$ to $6.2\,\%$, because magnesium took $86.8\,\%$ of the
charge instead of $32.0\,\%$ and the lithium simply had no current left to ride on. The cost
of a mole of lithium rose from $92$ to $475\unit{W\,h}$ for the same reason: the energy
denominator in \eqref{eq:energy-ae} shrank by a factor of nearly six while the energy itself
fell only a tenth.

Permselectivity tells the more careful story. It fell too, from $2.4982$ to $0.70745$, but
far less than recovery did --- a factor of $3.5$ against a factor of $6.2$. Dividing by the
feed ratio in \eqref{eq:perm-ae} takes out most, though plainly not all, of the four-fold
magnesium increase. What is left is a real loss of selectivity under a magnesium-rich feed,
not an artefact of the feed being different, and that is exactly the separation
\eqref{eq:perm-ae} exists to make.

One detail of AG is worth following, because it shows the checks and the calculations
working together. Minute 225 of AG's power log reads $100.06\unit{V}$ on the supply against
$10.0\unit{V}$ on the multimeter. \texttt{PairedValuesTooFarApart} clears the supply reading
and writes it to the journal; had it not, AG's highest voltage would read $100.06\unit{V}$
and its energy would be overstated. With the reading gone, the energy trapezoid bridges
minute 210 to minute 240 in one $30$-minute step worth $24.76\unit{W\,min}$, and the highest
voltage falls back to the $13.9\unit{V}$ at minute 0, which is what the table reports.

The balance gaps are diagnostics only; nothing is ever corrected by them. AE's lithium gap is
negative, meaning the concentrate and rinse together ended with $59.6\unit{mg}$ more lithium
than the feed appears to have lost, which is sampling and weighing error rather than lithium
appearing from nowhere. The magnesium gaps of $35\,\%$ and $45\,\%$ are mostly liquid left
behind in the sample bottles, and the rinse is counted as zero whenever it was never weighed,
which can only make a gap look worse --- the honest direction for a diagnostic to be wrong in.

\subsection{What the test suite pins}

The numbers above are AE's and AG's, and they are not literals anywhere in the repository ---
re-run \texttt{uv run limodel analyze} and the report recomputes them from the four CSV files
every time. What the suite pins is the arithmetic that produced them, one test file per
quantity under \texttt{tests/\allowbreak analyzer/\allowbreak calculate/\allowbreak
quantities/}, built on fixture runs whose
tares are round and whose tanks hold exactly $1.000$ or $2.000\unit{L}$ so every expected
milligram can be checked in the head.

So each choice this section leaned on has a test standing behind it. A fixture holding
$10\unit{mmol}$ each of lithium and magnesium asserts $t_{\mathrm{Li}} = 1/3$, which fails the
moment the charge weighting in \eqref{eq:tn} is dropped. Another asserts recovery of
$30\,\%$ and states in its own name that it measures against the feed's first minute and not
its last, so moving the denominator in \eqref{eq:recovery-ae} breaks the build. A third feeds the
same crossing from two different feeds and asserts permselectivity of $0.5$ and then $1.0$,
which is \eqref{eq:perm-ae} refusing to be rewritten without the feed ratio. The trapezoid in
\eqref{eq:charge-ae} is pinned by a steady $0.1\unit{A}$ for an hour reading $360\unit{C}$, and
by a log handed over out of minute order that must still read $900\unit{C}$.

A change that quietly altered any step of AE's chain would have to get past those first.

\section{Fitting the polynomial}

The modelling stage takes the \texttt{Result} rows the calculator handed on and makes, for each
membrane separately, three fitted equations and one search. Nothing here looks at a reading or
a sample again. The whole stage is \texttt{src/limodel/analyzer/model/}, and
\texttt{Surrogate.find\_models} is the only entry point.

Two things are decided before any arithmetic. First, the membrane is not a variable: the runs
are grouped by \texttt{cem} and the equation is fitted once per membrane.
\texttt{docs/polynomial.md} allowed either that or a 0/1 indicator; the code took the
per-membrane branch, which means the fits never borrow strength from each other. Second, a
membrane needs at least \texttt{Surrogate.FEWEST\_RUNS}, which is 3, because an intercept and
one slope is two coefficients and leave-one-out needs a row to spare.

On the 16 calculated runs that leaves two membranes fitted and two refused:

\begin{table}[htbp]
  \centering
  \caption{The four membranes in the current data, and what the fitting did with each.}
  \label{tab:membranes-fitted}
  \begin{tabular}{@{}lrl@{}}
    \toprule
    Membrane & Runs & Outcome \\
    \midrule
    GO-10            & 9 & fitted in both settings \\
    Pristine         & 5 & fitted in current density; feed ratio dropped \\
    PDA/PEI/15CE/BT  & 1 & not fitted: 1 run against the 3 a fit needs \\
    PDA-GO10/PEI/BT  & 1 & not fitted: 1 run against the 3 a fit needs \\
    \bottomrule
  \end{tabular}
\end{table}

A membrane that cannot be fitted is not silently skipped. It goes into
\texttt{SurrogateModels\allowbreak{}.not\_fitted} with the reason attached, and the report prints
it, because
``this membrane has one run'' is the next experiment worth booking rather than a footnote.

\subsection{The two design columns, logged and coded}

\texttt{DesignPoints} owns both settings of every run and the two translations between real
units and the units the fitting works in. The columns are
\texttt{current\_density\_ma\_cm2} and \texttt{feed\_mg\_li\_ratio} --- current density $J$ in
$\unit{mA/cm^2}$, and the feed's Mg/Li in $\unit{mol/mol}$. Both come from the calculate stage:
$J$ is the median supply current over the membrane area, and the feed ratio is worked out from
the run's own feed concentrations.

Each setting is logged first, then centred and scaled:

\begin{align}
  \label{eq:coding}
  x_1 &= \frac{\ln J - c_J}{h_J}, &
  c_J &= \tfrac{1}{2}\bigl(\ln J_{\max} + \ln J_{\min}\bigr), &
  h_J &= \tfrac{1}{2}\bigl(\ln J_{\max} - \ln J_{\min}\bigr), \\
  \label{eq:coding-r}
  x_2 &= \frac{\ln R - c_R}{h_R}, &
  c_R &= \tfrac{1}{2}\bigl(\ln R_{\max} + \ln R_{\min}\bigr), &
  h_R &= \tfrac{1}{2}\bigl(\ln R_{\max} - \ln R_{\min}\bigr).
\end{align}

The extremes in \eqref{eq:coding} and \eqref{eq:coding-r} are taken over the runs of the
membrane being fitted, so $x_1$ and $x_2$ run from $-1$ to $+1$ by construction. That is the
whole point of the coding, and the search later leans on it.

Two reasons for the two transforms, and they are separate reasons.

The logarithm is because each setting spans several fold rather than a few units. Across the
GO-10 runs $J$ goes from $4.50$ to $25.67\unit{mA/cm^2}$, a factor of $5.70$, and the feed ratio
from $1.43$ to $11.42\unit{mol/mol}$, a factor of $8.00$. On a raw scale a quadratic over that
span is dominated by whatever happens at the top end; on a log scale the steps between the
levels run are roughly even, so every run gets comparable say in the curvature.

The centring and scaling is because the squared terms would otherwise swamp the linear ones.
Logged but uncoded, the GO-10 current density column runs over $1.504$ to $3.245$, so its
square would run over $2.26$ to $10.53$ --- a column an order of magnitude larger than its own
linear partner, strongly correlated with it, and pushing the least-squares solve towards a
badly conditioned matrix. In coded units $x$ and $x^2$
both sit inside $[-1,1]$ and the coefficients are directly comparable: $b_1$ and $b_3$ are both
``change in the outcome across half the measured span''.

Table~\ref{tab:coded} is the coding for the runs as they stand.

\begin{table}[htbp]
  \centering
  \caption{The GO-10 design, in real units and coded. Centres and half-widths are
           $c_J = 2.374635$, $h_J = 0.870558$, $c_R = 1.395778$, $h_R = 1.039721$ in logged
           units; $e^{c_J} = 10.75\unit{mA/cm^2}$ and $e^{c_R} = 4.04\unit{mol/mol}$ are the
           geometric centres.}
  \label{tab:coded}
  \begin{tabular}{@{}lrrrr@{}}
    \toprule
    Run & $J$ ($\unit{mA/cm^2}$) & $R$ ($\unit{mol/mol}$) & $x_1$ & $x_2$ \\
    \midrule
    AE & 13.67 &  1.43 & $+0.276058$ & $-1.000000$ \\
    AF & 13.67 &  2.86 & $+0.276058$ & $-0.333333$ \\
    AG & 13.67 &  5.71 & $+0.276058$ & $+0.333333$ \\
    AH &  7.50 &  1.43 & $-0.413220$ & $-1.000000$ \\
    AI &  7.50 &  1.43 & $-0.413220$ & $-1.000000$ \\
    AJ &  4.50 &  1.43 & $-1.000000$ & $-1.000000$ \\
    AN & 25.67 &  1.43 & $+1.000000$ & $-1.000000$ \\
    AP & 25.67 &  1.43 & $+1.000000$ & $-1.000000$ \\
    AT & 13.67 & 11.42 & $+0.276058$ & $+1.000000$ \\
    \bottomrule
  \end{tabular}
\end{table}

The five Pristine runs --- AK, AL, AM, AO, AQ --- sit at $J = 13.67$, $7.50$, $4.50$, $25.67$
and $25.67\unit{mA/cm^2}$, all at $R = 1.43\unit{mol/mol}$. Their $x_1$ values are the same four
numbers as GO-10's, because the extremes of $J$ happen to be the same, and their $x_2$ is $0.00$
throughout.

The shape of Table~\ref{tab:coded} is worth a moment, because it governs everything after it.
Current density was moved only at the lowest feed ratio, and the feed ratio was moved only at
$13.67\unit{mA/cm^2}$. The design is a cross, not a grid.

The code's columns differ in two places from the table in the agreed
document, and neither difference is a code fault:

\begin{itemize}[nosep]
  \item The doc lists $R$ levels of $5, 10, 20, 40, 60$. Those are mass ratios --- the feed
        holds $400\unit{mg/L}$ Li against $2000$, $4000$, $8000$ or $16000\unit{mg/L}$ Mg. The
        code's column is $\unit{mol/mol}$, so the same four levels read as $1.43$, $2.86$,
        $5.71$ and $11.42$, a fixed factor of $M_{\mathrm{Li}}/M_{\mathrm{Mg}} = 0.285538$ apart.
        Because \eqref{eq:coding-r} logs then centres, a constant factor moves $c_R$ and nothing
        else --- the coded values are identical either way. The $60$ level has not been run.
  \item The doc lists $J$ levels of $3, 6, 12, 24$. Those are nominal. The code uses the median
        supply current actually logged, which gives $4.50$, $7.50$, $13.67$ and
        $25.67\unit{mA/cm^2}$. The fit uses what the supply held, not what it was asked for.
\end{itemize}

\subsection{A setting held at one value is dropped, not a refusal}

The feed ratio never moved across the Pristine runs. A naive fit would then have a column of
identical numbers, which after centring is a column of zeroes --- the solve would be rank
deficient and the whole membrane would be thrown out.

\texttt{DesignPoints} handles it by name instead. \texttt{find\_scales} gives that setting a
half-width of exactly $0$; \texttt{find\_kept\_settings} returns only the settings whose
half-width is non-zero; \texttt{find\_usable\_coded\_names} turns those into the coded names the
polynomial knows, and \texttt{find\_dropped\_settings} names the rest so the report can say so.
A dropped setting is coded as $0.0$, the one place it was ever run, and its terms are left out by
name, so the zero column never reaches a least-squares call.

Today this applies to Pristine and only Pristine. What it changes:

\begin{itemize}[nosep]
  \item The candidate terms shrink from four to two. A term is kept only if the setting
        under it varies, which \texttt{find\_candidates\_on} decides, so Pristine's
        candidates are \texttt{x1} and \texttt{x1\_squared}.
  \item The report says it out loud. The ``fitted in'' line reads \emph{current density, with
        the feed ratio held at one value across these runs and never tested}, which stops a fit
        in one setting being read as a fit in two.
  \item The search still returns a feed ratio. \texttt{find\_real\_from\_coded} multiplies the
        coded value by a half-width of zero, so whatever $x_2$ the grid offers reads back as
        $1.43\unit{mol/mol}$ --- the one value it was held at.
\end{itemize}

Both settings held at one value is a different matter: \texttt{find\_scales} then returns
\texttt{None}, and \texttt{Surrogate.find\_refusal\_for} refuses the membrane with \emph{no
setting that varies across its runs, so there is nothing to fit}. Runs that moved nothing carry
no information about anything.

\subsection{The candidate terms, and why there is no interaction}

\texttt{SecondOrderPolynomial.TERMS\_IN\_ORDER} is five terms, not six:

\begin{equation}
  \label{eq:polynomial}
  \hat{y} \;=\; b_0 \;+\; b_1 x_1 \;+\; b_2 x_2 \;+\; b_3 x_1^2 \;+\; b_4 x_2^2
\end{equation}

\texttt{docs/polynomial.md} writes a sixth term, $b_5 \, J \cdot R$, and calls it the term that
says whether the best current depends on how much Mg is in the feed. It is not in the code, and
the reason is arithmetic rather than preference: over these runs that column does not carry any
information the other three do not already carry.

Here is the proof rather than the assertion. Look again at Table~\ref{tab:coded}. Every run with
$x_2 \neq -1$ sits at $x_1 = 0.276058$, and every run with $x_1 \neq 0.276058$ sits at
$x_2 = -1$. Write $u = 0.276058056371$, the coded value of $13.67\unit{mA/cm^2}$. Then for every
one of the nine GO-10 rows,

\begin{equation}
  \label{eq:identity}
  x_1 x_2 \;=\; u \;-\; x_1 \;+\; u\,x_2 .
\end{equation}

Check it both ways. On the arm where $x_2 = -1$, the right side is $u - x_1 - u = -x_1$, which is
$x_1 x_2$. On the arm where $x_1 = u$, it is $u - u + u x_2 = u x_2$, which is again $x_1 x_2$.
Every row lies on one arm or the other, so \eqref{eq:identity} holds everywhere.

Regressing the $x_1 x_2$ column on $[\mathbf{1}, x_1, x_2]$ with \texttt{numpy.linalg.lstsq}
returns the coefficients $(+0.276058, -1.000000, +0.276058)$ and a largest absolute residual of
$6.7 \times 10^{-16}$ --- floating-point dust, not a near miss. Checking \eqref{eq:identity}
directly gives a largest error of $1.1 \times 10^{-16}$.

So the interaction is an exact linear combination of the intercept and the two linear terms, and
the matrix ranks say the same thing:

\begin{table}[htbp]
  \centering
  \caption{Rank of the design matrix over the fitted runs, from
           \texttt{numpy.linalg.matrix\_rank}.}
  \label{tab:rank}
  \begin{tabular}{@{}lrrr@{}}
    \toprule
    Columns & Rows & Columns & Rank \\
    \midrule
    \multicolumn{4}{@{}l}{\textit{GO-10, 9 runs}} \\
    $[\mathbf{1},\,x_1,\,x_2,\,x_1 x_2]$                     & 9 & 4 & 3 \\
    $[\mathbf{1},\,x_1,\,x_2,\,x_1^2,\,x_2^2]$               & 9 & 5 & 5 \\
    $[\mathbf{1},\,x_1,\,x_2,\,x_1^2,\,x_2^2,\,x_1 x_2]$     & 9 & 6 & 5 \\
    \midrule
    \multicolumn{4}{@{}l}{\textit{Pristine, 5 runs}} \\
    $[\mathbf{1},\,x_1,\,x_2,\,x_1 x_2]$                     & 5 & 4 & 2 \\
    $[\mathbf{1},\,x_1,\,x_2,\,x_1^2,\,x_2^2,\,x_1 x_2]$     & 5 & 6 & 3 \\
    \bottomrule
  \end{tabular}
\end{table}

The singular values of the full six-column GO-10 matrix make the same point:

\begin{equation}
  \label{eq:singular}
  \sigma \;=\; \bigl(4.6488,\; 2.6124,\; 1.8651,\; 0.9843,\; 0.7760,\; 0.0000\bigr)
\end{equation}

The last entry of \eqref{eq:singular} is zero, not small. A zero singular value means the
interaction coefficient does not exist: infinitely many values of $b_5$, each with matching
changes in $b_0$, $b_1$ and $b_2$, fit the nine rows exactly as well. Least squares would hand
back one of them, and it would mean nothing.

The five-term matrix in \eqref{eq:polynomial} is full rank, 5 of 5, so the squared terms are not
affected by this. It is the product that has no support.

Pristine is simpler still: $x_2 = 0$ throughout, so $x_1 x_2$ is literally a column of zeroes.

The code does not just omit the term --- it also refuses any fit that turns out rank deficient.
\texttt{SecondOrderPolynomial.find\_fit} reads the rank back from \texttt{lstsq} and returns
\texttt{None} when it is short of the number of terms, so a term set these rows cannot settle
can never be reported as if it had been settled.

Getting the interaction back needs runs, not code: one run at a high feed ratio and a different
current density would break \eqref{eq:identity} on its own. AT moved the feed ratio to
$11.42\unit{mol/mol}$ but stayed at $13.67\unit{mA/cm^2}$, so it sits on the same arm.

\subsection{Forward stepwise selection}

The agreed document, \texttt{docs/polynomial.md}, writes the full equation and fits it. The code
does not. It grows the equation one term at a time, in
\texttt{FittedOutcome.find\_best\_polynomial}:

\begin{enumerate}[nosep]
  \item Start with the intercept alone and record its leave-one-out error.
  \item Of the terms it is allowed to add, try each one and keep the one that lowers the error
        most.
  \item Stop when no remaining term lowers it.
\end{enumerate}

Two rules shape the search. A squared term waits for its own linear term:
\texttt{find\_terms\_it\_may\_add} will not offer \texttt{x2\_squared} until \texttt{x2} is in,
because a curvature with no slope under it cannot be read as anything. And ties go to the
simpler term --- candidates are scanned in the order \texttt{x1}, \texttt{x2},
\texttt{x1\_squared}, \texttt{x2\_squared}, and an error that merely equals the incumbent does
not displace it, nor does it keep the search going.

This is a departure from the agreed document, and the reason is that the runs cannot settle that
many coefficients. Nine GO-10 rows for a five-term equation leaves four degrees of freedom, and
two of the nine rows repeat two others --- so the design holds seven distinct points. Fitting
five coefficients to seven points produces an equation that passes close to everything and
predicts nothing. Pristine is worse: five runs, four distinct currents. Growing
the equation only while each new term pays for itself out of sample is how the fit is kept to the
size the data can support.

The honest caveat is that the reported error is the same quantity the search minimised, so it is
mildly optimistic as an estimate of error on a fresh run. The comparison in the next subsection
is what guards against reading too much into it.

Once the term set is chosen, the coefficients that get printed come from a final fit on
\emph{all} the rows: \texttt{find\_fit} refits the chosen terms over every row after the
search finishes. The leave-one-out folds are only ever used to score term sets, never to
produce the published coefficients.

Table~\ref{tab:stepwise} is the search as it actually ran, for all six fits.

\begin{table}[htbp]
  \centering
  \caption{Every step of the forward search, with the leave-one-out error each candidate
           reached. The chosen term at each step is marked; a dash means the term was not
           offered, either because it was already in or because its linear partner was not.}
  \label{tab:stepwise}
  \begin{tabular}{@{}lrrrr@{}}
    \toprule
    Fit so far & \texttt{x1} & \texttt{x2} & \texttt{x1\_squared} & \texttt{x2\_squared} \\
    \midrule
    \multicolumn{5}{@{}l}{\textit{GO-10, permselectivity --- intercept alone: 0.7678}} \\
    intercept                     &  0.8683 & \textbf{0.5611} & ---    & ---    \\
    intercept, \texttt{x2}        &  0.7584 & ---             & ---    & 0.6051 \\
    \midrule
    \multicolumn{5}{@{}l}{\textit{GO-10, recovery \% --- intercept alone: 15.0521}} \\
    intercept                     & 16.2443 & \textbf{7.8979} & ---    & ---    \\
    intercept, \texttt{x2}        &  9.9497 & ---             & ---    & \textbf{7.1218} \\
    $\dots$, \texttt{x2\_squared} &  9.4710 & ---             & ---    & ---    \\
    \midrule
    \multicolumn{5}{@{}l}{\textit{GO-10, volts --- intercept alone: 4.0593}} \\
    intercept                     & \textbf{1.3385} & 4.3523  & ---    & ---    \\
    intercept, \texttt{x1}        & ---     & 1.5419          & 1.4018 & ---    \\
    \midrule
    \multicolumn{5}{@{}l}{\textit{Pristine, permselectivity --- intercept alone: 0.0614}} \\
    intercept                     &  0.0741 & ---             & ---    & ---    \\
    \midrule
    \multicolumn{5}{@{}l}{\textit{Pristine, recovery \% --- intercept alone: 2.9138}} \\
    intercept                     &  3.6629 & ---             & ---    & ---    \\
    \midrule
    \multicolumn{5}{@{}l}{\textit{Pristine, volts --- intercept alone: 5.8947}} \\
    intercept                     & \textbf{1.0682} & ---     & ---    & ---    \\
    intercept, \texttt{x1}        & ---     & ---             & 2.9568 & ---    \\
    \bottomrule
  \end{tabular}
\end{table}

Three things in that table are worth reading twice. No fit ever reached three terms except
GO-10's recovery, which took \texttt{x2} and then \texttt{x2\_squared}. Current density never
earned a place in a permselectivity or recovery fit --- for GO-10 it made both errors
\emph{worse} than the intercept alone. And two Pristine fits kept nothing at all: the best
prediction of its permselectivity or recovery at any setting is the average of the five runs.

\subsection{Leave-one-out, not a hold-out set}

\texttt{LeaveOneOutError.find\_rmse} is the only error measure in the stage. Each row takes its
turn being held back; the polynomial is fitted on the rest and asked to predict the one it never
saw; the misses are squared, averaged and rooted:

\begin{equation}
  \label{eq:loo}
  \mathrm{RMSE}_{\text{LOO}} \;=\; \sqrt{\frac{1}{n}\sum_{i=1}^{n}
    \bigl(\hat{y}^{(-i)}(x_i) - y_i\bigr)^2}
\end{equation}

Every row in \eqref{eq:loo} is predicted by a fit that did not include it, so the number is out
of sample $n$ times over. It returns \texttt{None} rather than a number whenever a fold cannot
be fitted, whenever an outcome is missing, and whenever there are fewer than two rows --- and the
stepwise search treats \texttt{None} as ``this term cannot be judged'' and moves on, which is
why a term set the data cannot settle never wins by default.

A train/test split was rejected for one reason: with nine rows a hold-out set measures nothing.
Hold two GO-10 runs back and the error is an average of two squared misses; swap which two and
the number moves more than the differences between the models it is meant to be choosing
between. Worse, the seven rows left would not span the design --- hold back AT and the feed
ratio no longer reaches $11.42\unit{mol/mol}$, so the model being tested is not the model being
proposed. Leave-one-out spends no rows permanently and produces one honest out-of-sample miss
per run instead of two in total.

The real test set is arriving on its own. AV, AW, AX and AY have no lab results yet, so they are
dropped today; when their results land they are unseen runs that no fit has touched.

\subsection{Replicate scatter: the yardstick}

An error of $0.561$ on permselectivity is meaningless on its own. It needs something to be
compared against, and the only fair comparison is how far apart the bench's own repeats came
out. \texttt{ReplicateScatter} works that out.

A repeat is two runs at the same membrane and the same two settings, matched to two decimal
places --- \texttt{ReplicateScatter.PLACES}. The rounding is there because the settings are
computed, not typed in: $25.666666666666668$ and $25.666666666666664$ are the same experiment,
and matching the raw floats would miss it. Runs missing the outcome or either setting are passed
over.

Within each group the spread is pooled across groups rather than averaged, which is the standard
way to combine two small samples of the same underlying noise:

\begin{equation}
  \label{eq:pooled}
  s \;=\; \sqrt{\frac{\sum_{g}\sum_{i \in g}\bigl(y_i - \bar{y}_g\bigr)^2}
                     {\sum_{g}\bigl(n_g - 1\bigr)}}
\end{equation}

The denominator in \eqref{eq:pooled} is the degrees of freedom, and it is reported beside the
spread so a reader can see how thin the yardstick is. No group with a repeat in it means no
yardstick at all, and \texttt{find\_spread} returns \texttt{None}.

Three design points repeat in the current data:

\begin{table}[htbp]
  \centering
  \caption{The replicated design points and what each pair gave, with the pooled spread from
           \eqref{eq:pooled}. $J$ is in $\unit{mA/cm^2}$ and $R$ in $\unit{mol/mol}$. GO-10
           pools two pairs; Pristine has one.}
  \label{tab:replicates}
  \begin{tabular}{@{}lrrlll@{}}
    \toprule
    Runs & $J$ & $R$ & Permselectivity & Recovery (\%) & Volts \\
    \midrule
    AH, AI &  7.50 & 1.43 & 2.3881, 2.0322 & 36.59, 37.66 & 10.55, 10.22 \\
    AN, AP & 25.67 & 1.43 & 1.9177, 1.3217 & 39.50, 26.32 & 18.54, 16.91 \\
    \multicolumn{3}{@{}l}{GO-10 pooled, 2 d.o.f.}
           & 0.3471 & 6.6124 & 0.8315 \\
    \midrule
    AO, AQ & 25.67 & 1.43 & 0.2115, 0.2986 &  7.83, 11.38 & 16.84, 17.04 \\
    \multicolumn{3}{@{}l}{Pristine pooled, 1 d.o.f.}
           & 0.0615 & 2.5075 & 0.1414 \\
    \bottomrule
  \end{tabular}
\end{table}

Look at the GO-10 recovery pair AN and AP: $39.50\%$ and $26.32\%$ at the same setting on the
same membrane. That single pair is most of the $6.6124$ pooled spread, and it is the honest
statement of how reproducible recovery is at $25.67\unit{mA/cm^2}$ today.

This is why \texttt{FittedOutcome} carries the spread alongside the coefficients and the error:
the three travel together because an equation on its own cannot be judged. A fit that misses by
more than the bench's own repeats has learnt nothing the repeats did not already say --- it has
fitted the scatter. The report makes that call outright rather than printing two numbers and
leaving it to the reader, and it makes it at the three decimals it prints, so two numbers that
display the same are not treated as separated.

\subsection{The gated search for the best setup}

The question is not ``what is the most selective setting'' but ``what is the most selective
setting that is still worth running''. So all three fitted outcomes are needed:
permselectivity to maximise, and recovery and voltage to gate.
\texttt{BestSetupUnderTheGates} does this per membrane.

\begin{itemize}[nosep]
  \item \textbf{The grid.} \texttt{STEPS} is 21 levels per setting, evenly spaced in coded units
        from $-1$ to $+1$, which is $0.1$ apart --- $441$ combinations in all. Because the coding
        is on logs, those levels are geometric in real units: each step multiplies $J$ by
        $1.0910$ and the feed ratio by $1.1096$.
  \item \textbf{The sweep stays inside the measured range.} Coded $-1$ and $+1$ are exactly the
        lowest and highest run of the membrane being searched, so sweeping that span and no
        further is what keeps the answer to a setting the bench has evidence for. For GO-10 that
        is $4.50$ to $25.67\unit{mA/cm^2}$ and $1.43$ to $11.42\unit{mol/mol}$. A quadratic
        extrapolated past its last data point curves away fast and in whichever direction the
        algebra happens to point; the recommendation would be a reading taken past the last run.
  \item \textbf{The recovery floor.} \texttt{RECOVERY\_FLOOR\_PCT} is $35.0$. Less lithium than
        that recovered is not worth a run.
  \item \textbf{The voltage ceiling.} \texttt{VOLTAGE\_CEILING\_V} is $20.0$. The supply and the
        stack are not asked to go above it.
  \item \textbf{The winner.} Of the points that clear both gates, the one with the highest
        predicted permselectivity, converted back to $\unit{mA/cm^2}$ and $\unit{mol/mol}$ by
        \texttt{find\_real\_from\_coded}.
\end{itemize}

Both gates must be clear; \texttt{find\_gates\_failed\_at} returns the list of gates a point
falls foul of, and only an empty list survives. A point where any one of the three outcomes
cannot be predicted returns \texttt{None} and is counted as neither passed nor failed.

Refusing every point is a finding and not an error, and the code is built around that. It keeps a
count of how many of the $441$ points each gate turned away, and \texttt{find\_refusal} names the
gate responsible in a sentence the report prints. ``Nothing in the range you have tested is worth
running'' is a true and useful answer to the question that was asked.

That is exactly what happens today. For both fitted membranes, all $441$ points fail the recovery
floor and $0$ fail the voltage ceiling:

\begin{table}[htbp]
  \centering
  \caption{The gated search on the current fits. The highest predicted values anywhere on the
           grid, against the two limits.}
  \label{tab:gates}
  \begin{tabular}{@{}lrrrrr@{}}
    \toprule
    Membrane & Points & Best recovery (\%) & Floor (\%) & Worst volts & Ceiling \\
    \midrule
    GO-10    & 441 & 33.285 & 35.0 & 18.246 & 20.0 \\
    Pristine & 441 &  9.057 & 35.0 & 17.320 & 20.0 \\
    \bottomrule
  \end{tabular}
\end{table}

GO-10 comes close --- $33.285\%$ against a floor of $35\%$ --- and the voltage was never the
binding constraint for either membrane. Pristine fails by a mile for a simple reason: its
recovery fit kept no terms, so it predicts $9.057\%$ everywhere, which is the mean of its five
runs.

\subsection{What the fits are today}

Table~\ref{tab:fits} is what \texttt{uv run limodel analyze} prints. The terms column is the
equation exactly as the report writes it; the error is the leave-one-out RMSE from
\eqref{eq:loo}; the scatter is the pooled replicate spread from \eqref{eq:pooled} with its
degrees of freedom.

\begin{table}[htbp]
  \centering
  \caption{Every fitted equation in the current data, with its leave-one-out error and the
           pooled bench scatter beside it; the scatter's degrees of freedom are in brackets.
           Only Pristine's permselectivity is inside the scatter.}
  \label{tab:fits}
  \begin{tabular}{@{}llrr@{}}
    \toprule
    Outcome & Terms as printed & Error & Scatter \\
    \midrule
    \multicolumn{4}{@{}l}{\textit{GO-10 --- fitted in current density and the feed ratio}} \\
    permselectivity & \texttt{intercept +1.126, x2 -0.712}      & 0.561 & 0.347 (2) \\
    recovery \%     & \texttt{intercept +11.58, x2 -15.63,}
                                                                 & 7.122 & 6.612 (2) \\
                    & \texttt{x2\_squared +6.08}                &       &           \\
    volts           & \texttt{intercept +12.7, x1 +5.545}       & 1.339 & 0.832 (2) \\
    \midrule
    \multicolumn{4}{@{}l}{\textit{Pristine --- fitted in current density, with the feed ratio
                                   held at one value}} \\
    permselectivity & \texttt{intercept +0.2326}                & 0.061 & 0.062 (1) \\
    recovery \%     & \texttt{intercept +9.057}                 & 2.914 & 2.508 (1) \\
    volts           & \texttt{intercept +11.41, x1 +5.905}      & 1.068 & 0.141 (1) \\
    \bottomrule
  \end{tabular}
\end{table}

Five of the six fits print the sentence \emph{which is worse than the bench repeating itself, so
this fit has earned no conclusion}. Neither membrane gets a best setup; both print \emph{nothing
inside the measured range cleared the lithium recovery floor}.

Read plainly, that is the state of the modelling. One fit --- Pristine permselectivity, at
$0.061$ against a scatter of $0.062$ --- is inside the bench's own noise, and it is the fit that
kept no terms at all, so what it says is that Pristine's permselectivity is about $0.23$
regardless of current density. Everything else misses the held-out run by more than two repeats
of the same run differ from each other.

The two coefficients that do point somewhere are worth naming even so, because they are the only
two the search kept more than once. GO-10's permselectivity falls with the feed ratio
($-0.712$ across half the coded span, roughly a factor of $2.83$ in Mg/Li), and voltage rises
with current density on both membranes by about the same amount ($+5.545$ and $+5.905$ volts
across half the span). Neither is a surprise, and neither has cleared the yardstick.

The useful conclusion is about the design rather than the chemistry. The cross in
Table~\ref{tab:coded} is why there is no interaction term, five of six errors are above the
scatter, and the two thinnest membranes have one run each. Runs at a high feed ratio and a
current density other than $13.67\unit{mA/cm^2}$ would break \eqref{eq:identity} and let the
interaction be fitted; a third run at $25.67\unit{mA/cm^2}$ on GO-10 would tighten the recovery
yardstick that AN and AP currently set at $6.6$ percentage points.

\section{What these runs can and cannot answer}
\label{sec:limits}

Everything in this section comes from the sixteen runs the analyzer used on the last
run of \texttt{limodel analyze}, and every number was read off the code rather than
recalled. Nothing here is a criticism of the code. The code is honest about what it
was handed; this section is about what was handed to it.

Twenty-one runs are in \texttt{data/runs.csv}. Sixteen reach the fitting. The five that
do not are dealt with in \S\ref{sec:limits:waiting}.

\subsection{The design is a cross, not a grid}
\label{sec:limits:cross}

A second-order polynomial in two settings wants a grid: every current density crossed
with every feed ratio. These runs are not that. They are two arms meeting at one corner.

The first arm moves current density with the feed ratio pinned at its lowest value,
$R = 1.43\unit{mol/mol}$. Table~\ref{tab:limits:current-arm} is every run on it.

\begin{table}[htbp]
  \centering
  \caption{The current-density arm. Feed ratio $1.43\unit{mol/mol}$ on all thirteen runs.
    Current density is the median supply current of the run over the membrane area, which
    is why the levels are not round numbers.}
  \label{tab:limits:current-arm}
  \begin{tabular}{l l l}
    \toprule
    $J$ ($\unit{mA/cm^2}$) & Membrane & Runs \\
    \midrule
    4.50  & GO-10            & AJ \\
          & Pristine         & AM \\
    7.50  & GO-10            & AH, AI \\
          & Pristine         & AL \\
          & PDA/PEI/15CE/BT  & AR \\
          & PDA-GO10/PEI/BT  & AS \\
    13.67 & GO-10            & AE \\
          & Pristine         & AK \\
    25.67 & GO-10            & AN, AP \\
          & Pristine         & AO, AQ \\
    \bottomrule
  \end{tabular}
\end{table}

The second arm moves the feed ratio with the current density pinned at
$13.67\unit{mA/cm^2}$, and it exists on one membrane only.

\begin{table}[htbp]
  \centering
  \caption{The feed-ratio arm. Current density $13.67\unit{mA/cm^2}$, membrane GO-10, on
    all four runs. AE is the corner the two arms share.}
  \label{tab:limits:ratio-arm}
  \begin{tabular}{l l l}
    \toprule
    $R$ ($\unit{mol/mol}$) & Membrane & Run \\
    \midrule
    1.43  & GO-10 & AE \\
    2.86  & GO-10 & AF \\
    5.71  & GO-10 & AG \\
    11.42 & GO-10 & AT \\
    \bottomrule
  \end{tabular}
\end{table}

Put the two tables side by side and the shape is plain. Four current-density levels exist
and four feed-ratio levels exist, so a full grid would hold sixteen cells. Seven are
occupied. Every run off the corner sits on one arm or the other, never in the middle of
the field.

This is why \texttt{SecondOrderPolynomial} has no $x_1x_2$ term. The code says so in a
comment beside \texttt{TERMS\_IN\_ORDER}, and \texttt{CANDIDATE\_TERMS} in
\texttt{fitted\_outcome.py} never offers one. An interaction term answers the question
\emph{does the best current density depend on how much magnesium is in the feed}. To
answer it, the feed ratio has to move at more than one current density. In these runs it
moves at one. The term is not missing from the code because it was forgotten. It is
missing because there is nothing in the data for it to be fitted to. That is a fact about
the experiments, and no change to the code can repair it.

There is a second thing the tables do not show, and it matters more than the missing
interaction. The four runs at $25.67\unit{mA/cm^2}$ --- AN, AO, AP and AQ --- all started
their concentrate at $5\unit{g/L}$. Every other run started it at $0.5\unit{g/L}$. AN and
AO started from sodium chloride; AP and AQ from lithium chloride. So the top level of the
current-density arm changes the concentrate at the same time as it changes the current,
and \texttt{Result} carries neither \texttt{conc\_start\_salt} nor
\texttt{conc\_start\_g\_l}, so the fitting cannot see the change and cannot allow for it.
Any curvature the fit finds at the top of the current range is a current effect and a
concentrate effect added together, with no way to tell the two apart.

\subsection{What the bench gets when it repeats itself}
\label{sec:limits:replicates}

\texttt{ReplicateScatter} groups runs by membrane and by both settings rounded to two
places, then pools the scatter inside the groups. Three groups have more than one run.
Table~\ref{tab:limits:replicates} gives them with the gap between the two members.

\begin{table}[htbp]
  \centering
  \caption{Every repeated design point among the sixteen runs, and the gap between the
    pair in each of the three fitted outcomes. Feed ratio is $1.43\unit{mol/mol}$ on all
    six runs.}
  \label{tab:limits:replicates}
  \begin{tabular}{l l r r r r}
    \toprule
    & & $J$ & Recovery & & Voltage \\
    Pair & Membrane & ($\unit{mA/cm^2}$) & (points) & $P$ & ($\unit{V}$) \\
    \midrule
    AH, AI & GO-10    & 7.50  & 1.07  & 0.356 & 0.33 \\
    AN, AP & GO-10    & 25.67 & 13.18 & 0.596 & 1.63 \\
    AO, AQ & Pristine & 25.67 & 3.55  & 0.087 & 0.20 \\
    \bottomrule
  \end{tabular}
\end{table}

Pooled, GO-10's two repeats give a standard deviation of 6.61 points of recovery, 0.347
in permselectivity and $0.832\unit{V}$, each on two degrees of freedom. Pristine's one
repeat gives 2.51 points, 0.062 and $0.141\unit{V}$ on one degree of freedom. These are
the numbers the report prints as the bench scatter.

Now put the largest gap next to the effect the fit is chasing. Along the current-density
arm, GO-10's recovery reads 21.1\% at $4.50\unit{mA/cm^2}$, 36.6\% and 37.7\% at 7.50,
38.3\% at 13.67, and 39.5\% and 26.3\% at 25.67. Averaging the pairs, the step from 7.50
to 13.67 is 1.2 points and the step from 13.67 to 25.67 is $-5.4$ points. The AN--AP gap
is 13.2 points. It is larger than both of those steps put together, and larger than the
whole span of the level averages above $4.50\unit{mA/cm^2}$. Only the first step, 4.50 to
7.50 at 16.0 points, is bigger than the scatter at a single point.

Permselectivity is no better placed. The AN--AP gap of 0.596 is twice the 0.288 that
separates the 7.50 and 13.67 levels, and most of the 0.878 that separates 13.67 from
25.67. Voltage is the one outcome that survives the comparison: the arm climbs
$6.59 \to 10.39 \to 16.0 \to 17.73\unit{V}$ and every step is larger than the
$1.63\unit{V}$ gap.

So for two of the three outcomes the bench's own scatter at one setting is as large as the
difference between settings. The report says the same thing in the language of error:
GO-10's leave-one-out error is 0.561 against a scatter of 0.347 in permselectivity and
7.12 points against 6.61 in recovery, and in both cases it prints \emph{this fit has
earned no conclusion}. No amount of fitting fixes this. Until a repeat at a given setting
lands closer than the difference between settings, a curve drawn through these points
cannot be told apart from a curve drawn through the scatter.

One caveat about that 13.2-point gap, because it is the number doing most of the damage.
AN and AP are a repeat only in the two columns the grouping looks at. They differ in the
concentrate they started from, sodium chloride against lithium chloride, both at
$5\unit{g/L}$. The same is true of AO and AQ. AH and AI are the only pair in the data
whose set-up columns all agree. So part of what is being charged to bench scatter may be
a real effect of the starting concentrate --- which would be better news for the
repeatability and worse news for the design, since it means the top current level is
confounded twice over.

\subsection{Two membranes cannot be fitted at all}
\label{sec:limits:membranes}

Of the sixteen runs, nine are GO-10 and five are Pristine. The other two are a single
run each of PDA/PEI/15CE/BT (AR) and of PDA-GO10/PEI/BT (AS).
\texttt{Surrogate.FEWEST\_RUNS} is 3: an intercept and one
slope are two coefficients, and leave-one-out needs a row to spare. One run is not three,
so the report prints

\begin{quote}
\texttt{not fitted}\\
\texttt{PDA/PEI/15CE/BT has too few runs to fit,}\\
\texttt{1 against the 3 a fit needs}
\end{quote}

and the same for PDA-GO10/PEI/BT. Both of those runs sit at the same point, GO-10's
best-repeated setting of $7.50\unit{mA/cm^2}$ and $R = 1.43$. They give a permselectivity
of 1.43 and 1.23 against GO-10's 2.39 and 2.03 and Pristine's 0.28 at that same point, so
they are a reading, but one reading each. Nothing can be said about how either coating
responds to current density or to feed ratio, because neither was ever moved on them.

Pristine is fitted, but only in one direction. Its five runs never leave
$R = 1.43\unit{mol/mol}$, so \texttt{DesignPoints.find\_dropped\_settings} reports the
feed ratio as dropped and the report says so in words: \emph{fitted in current density,
with the feed ratio held at one value across these runs and never tested}. A Pristine
prediction at a higher feed ratio is not a cautious extrapolation. It is a number the
runs contain no information about at all.

\subsection{One polynomial over all sixteen runs}
\label{sec:limits:pooled}

The analyzer fits each membrane separately. The obvious alternative is one polynomial
over everything, which is what \texttt{docs/polynomial.md} describes. It is worth running,
because what comes back is the clearest statement of the problem.

The fit uses the project's own machinery: \texttt{DesignPoints} to code the settings and
\texttt{SecondOrderPolynomial} with all five of \texttt{TERMS\_IN\_ORDER}, forced in
rather than grown term by term. The coded settings are logged, then centred and scaled so
the lowest run sits at $-1$ and the highest at $+1$:

\begin{align}
  x_1 &= \frac{\ln\!\left(J / 10.747\right)}{0.87056}, \label{eq:limits:x1} \\
  x_2 &= \frac{\ln\!\left(R / 4.0381\right)}{1.03972}. \label{eq:limits:x2}
\end{align}

So $x_1 = 0$ is $10.75\unit{mA/cm^2}$, one coded unit is a factor of 2.39 in current
density, and one coded unit in $x_2$ is a factor of 2.83 in feed ratio. Fitting lithium
recovery in percent over all sixteen runs gives

\begin{equation}
  \widehat{\text{recovery}} = 12.09 + 2.03\,x_1 - 12.80\,x_2 - 8.97\,x_1^2 + 2.55\,x_2^2 .
  \label{eq:limits:pooled}
\end{equation}

Three numbers judge it:

\begin{itemize}[nosep]
  \item $R^2 = 0.299$. The surface accounts for 30\% of the variation in recovery across
        the sixteen runs.
  \item Leave-one-out error, through \texttt{LeaveOneOutError}, is 13.98 points.
  \item The standard deviation of the sixteen recoveries themselves is 12.96 points.
\end{itemize}

The leave-one-out error is the larger. That has a plain meaning. Hold a run back, fit
\eqref{eq:limits:pooled} to the other fifteen, predict the held-back run, and the
prediction misses by more than you would have missed by ignoring the settings entirely
and quoting the average recovery. The comparison can be made exactly: fitting an
intercept alone under the same leave-one-out loop gives 13.39 points. So the four extra
terms make out-of-sample prediction worse, not better. This is why
\texttt{FittedOutcome.find\_best\_polynomial} grows a model one term at a time and stops
when the next term does not lower the leave-one-out error --- a five-term fit forced onto
sixteen runs spread over seven design cells is fitting the scatter.

The surface's own peak says the same thing from another direction. Searched over the
coded box the runs cover, \eqref{eq:limits:pooled} tops out at 27.6\% recovery, at
$J = 11.9\unit{mA/cm^2}$ and $R = 1.43\unit{mol/mol}$. Four runs already measured beat it:
AE at 38.3\%, AH at 36.6\%, AI at 37.7\% and AN at 39.5\%, all GO-10 at $R = 1.43$. A
best-of-surface that sits more than ten points below four measurements in hand is not an
optimum. It is an average of membranes that behave nothing alike, and it is also why the
per-membrane report ends with \emph{no best setup: nothing inside the measured range
cleared the lithium recovery floor}. GO-10's own fitted recovery peaks at 33.3\%, under
the $35\%$ floor in \texttt{BestSetupUnderTheGates.RECOVERY\_FLOOR\_PCT}, while the same
four runs measured above it.

\texttt{docs/polynomial.md} allows the membrane in as a 0/1 indicator rather than as a
separate fit per membrane. Adding one such column, 1 for GO-10 and 0 for everything else,
to the same five terms gives

\begin{equation}
  \widehat{\text{recovery}} = -8.59 + 1.61\,x_1 - 18.13\,x_2 - 10.16\,x_1^2
    + 9.21\,x_2^2 + 20.15\,g ,
  \label{eq:limits:indicator}
\end{equation}

where $g = 1$ on a GO-10 run. $R^2$ rises from 0.299 to 0.819 and the leave-one-out error
falls from 13.98 to 7.65 points, now comfortably under the 12.96-point standard deviation.
One binary column explains more than all four setting terms together, and its coefficient
is 20 points of recovery --- larger than anything either setting does across its whole
measured range.

Read that carefully. It does not mean \eqref{eq:limits:indicator} is a good model. It
means the membrane is the dominant variable in this data set and the two dialled settings
are secondary to it. The indicator fit also still leans on the Pristine runs to pin
$x_2^2$, and Pristine never left $R = 1.43$, so the $+9.21$ on $x_2^2$ is carried by GO-10
alone. The per-membrane approach the analyzer actually takes is the more honest reading of
the same fact: if the membrane is worth 20 points, fit the membranes apart.

\subsection{The runs waiting on the lab}
\label{sec:limits:waiting}

Five of the twenty-one runs are dropped, all for the same reason, which the journal
records as \emph{no lab results at both ends}. Without a sample at the first and last
minute there is no amount that crossed, so there is no recovery and no permselectivity to
fit. Table~\ref{tab:limits:waiting} is the five, with the settings read from their own
rows and power logs.

\begin{table}[htbp]
  \centering
  \caption{The dropped runs and the settings they were made at. Current density is the
    median of each power log over the $6\unit{cm^2}$ area; all five come out at
    $13.5\unit{mA/cm^2}$.}
  \label{tab:limits:waiting}
  \begin{tabular}{l l r r}
    \toprule
    Run & Membrane & $J$ ($\unit{mA/cm^2}$) & $R$ ($\unit{mol/mol}$) \\
    \midrule
    AX & Pristine & 13.5 & 2.86 \\
    AY & Pristine & 13.5 & 5.71 \\
    AU & Pristine & 13.5 & 11.42 \\
    AV & Pristine & 13.5 & 17.13 \\
    AW & GO10     & 13.5 & 17.13 \\
    \bottomrule
  \end{tabular}
\end{table}

Four of the five are a Pristine feed-ratio sweep at one current density: 2.86, 5.71, 11.42
and 17.13, the same ladder GO-10 already has, with one rung added past its top. That is
exactly the dimension \S\ref{sec:limits:cross} found missing --- Pristine currently has no
feed-ratio movement at all, which is why the report calls the feed ratio dropped and
never tested on it. The fifth, AW, extends GO-10's own sweep to 17.13.

Their current density, $13.5\unit{mA/cm^2}$, is the same nominal setting as the existing
feed-ratio arm at $13.67\unit{mA/cm^2}$; the $0.17$ difference is the median of one power
log against another, not a different dial position. The concentrate start on all five is
lithium chloride at $0.5\unit{g/L}$, matching the existing arm, so these runs do not carry
the confound described in \S\ref{sec:limits:cross}.

Two things about them are worth knowing before the results land. First, the sweep is still
a sweep: it adds a second arm at the same current density, so after it arrives the feed
ratio will move on two membranes but still at one current density, and there will still be
no support for an interaction term. Second, AW's membrane is spelled \texttt{GO10} in
\texttt{data/runs.csv} while the nine fitted runs are spelled \texttt{GO-10}. The field is
a plain string --- \texttt{cem} in \texttt{models/run.py} and in the schema --- and
\texttt{Surrogate.find\_membranes} groups on it exactly as written, so AW will arrive as a
fourth membrane with one run and be refused for having too few runs, unless the spelling
is made to match.

\subsection{The balance gaps, and what they mean for quoting any number}
\label{sec:limits:gaps}

\texttt{BalanceGap} reports, as a percentage, how much of what the dilute lost was never
found in the concentrate and the rinse. It is a diagnostic. Nothing is ever corrected by
it. But it is the best single check on whether the lab numbers a run was built from hang
together, so it belongs in any judgement of how far these results can be pushed.

Over the sixteen runs the lithium gaps run from $-272.5\%$ to $+71.3\%$, with a median of
$+8.5\%$. Ten are positive, six negative. Seven of the sixteen are more than 25\% away
from zero in either direction, and only five are inside $\pm 10\%$. The magnesium gaps are
worse: $-76.8\%$ to $+67.6\%$, median $+17.5\%$, ten of sixteen outside $\pm 25\%$ and
only three inside $\pm 10\%$.

Both signs have a meaning, and only one of them is benign.

A positive gap means the dilute lost more than the other two tanks gained. That is what
hold-up looks like: liquid left in sample bottles and lines. \texttt{BalanceGap} counts an
unweighed rinse as zero gain, which it explains as the honest way for a diagnostic to be
wrong, and that choice pushes gaps positive. So a modest positive gap is roughly what a
sound run should show, and the median of $+8.5\%$ in lithium is consistent with that.

A negative gap means the concentrate and rinse together gained more metal than the dilute
lost. As mass, that cannot happen. Six runs show it in lithium and six in magnesium, so in
those runs something in the concentrations, the dilution factors or the tank weights does
not reconcile, and the amounts that crossed --- the numerator of both permselectivity and
recovery --- inherit that.

Two further points about size. The gap is a ratio, so it blows up when the denominator is
small. AT's $-272.5\%$ is the clearest case: its dilute lost $5.1\unit{mg}$ of lithium
over the whole run while the concentrate gained $18.8\unit{mg}$. A $14\unit{mg}$
discrepancy on a $5\unit{mg}$ loss is a huge percentage and a small absolute error, and AT
is the highest feed ratio in the data, the far end of the feed-ratio arm. The point at the
end of the lever is the one with the least lithium movement behind it. The other is that
the magnesium gaps are consistently larger than the lithium ones, which bears directly on
permselectivity, since $P$ divides a lithium transport number by a magnesium one and the
magnesium side is the noisier of the two.

The practical reading is narrow and worth stating plainly. A recovery or a permselectivity
quoted from one of these runs should be quoted next to that run's gaps. Five of sixteen
runs balance lithium inside $\pm 10\%$ and three balance magnesium that well. The report
prints both gap columns for every run for this reason, and a number taken out of that
table without them is being quoted more precisely than it was measured.


\end{document}
